\documentclass[12pt]{article}
\usepackage[utf8]{inputenc}
\usepackage[LGR,T1]{fontenc}
\usepackage[greek.ancient,english]{babel}


\usepackage[letterpaper, margin = 1in]{geometry}
\usepackage{fancyhdr}
\usepackage{titling}
\usepackage{titlesec}
\usepackage{hyperref}
\usepackage{lastpage}
\usepackage{lipsum}

\usepackage[nocenter]{qtree}
\usepackage{tree-dvips}
\usepackage{linguex}

\usepackage[backend=biber, style=apa]{biblatex}
\addbibresource{deverbals.bib}
\usepackage{csquotes}
\DeclareLanguageMapping{english}{english-apa}

\usepackage{soul}
\usepackage{tabularx}
\usepackage{hanging}
\usepackage{multirow}
\usepackage{array}
\usepackage{multirow}
\renewcommand\tabularxcolumn[1]{m{#1}}
\usepackage{amssymb}

\usepackage{gb4e}

\title{CSCL 1301W: Teaching Materials and Text Notes}
\author{Ian White}
\date{}

\makeatletter
\renewcommand{\maketitle}
{
\begin{flushleft}
\textbf{{\large\@title}}
\newline
\@author
\end{flushleft}
}
\makeatother

\pagestyle{fancy}
\fancyhf{}
\lhead{Ian Thomas White}
\chead{\small{CSCL1301W, Spring 2020}}
\rhead{\thepage \ of \pageref{LastPage}}

\widowpenalty=100000
\clubpenalty=100000

\titleformat*{\section}{\medium\bfseries}
\titleformat*{\subsection}{\small\bfseries}
\titleformat*{\subsubsection}{\small\bfseries}

\makeatletter
\renewcommand\part{%
  \par\bigskip
  \secdef\@part\@spart}
\def\@part[#1]#2{%
  \refstepcounter{part}%
  \addcontentsline{toc}{part}{#1}%
  {\noindent\large\bfseries #2\par}%
  \nobreak\smallskip}
\def\@spart#1{%
  {\noindent\large\bfseries #1\par}%
  \nobreak\smallskip}
\makeatother

\begin{document}

\thispagestyle{empty}

\begin{flushright}
\small{CSCL1301W: \textbf{Reading Culture, Spring 2020} \\ instructor: \textbf{Ian White} \\ contact: \textbf{whit2074@umn.edu}}
\end{flushright}

\renewcommand{\labelitemi}{$\square$}

\renewcommand{\labelenumi}{(\Roman{enumi})}
\renewcommand{\labelenumii}{(\Alph{enumii})}
\renewcommand{\labelenumiii}{(\roman{enumiii})}
\renewcommand{\labelenumiV}{(\alph{enumiV})}
\renewcommand{\labelenumV}{(\alph{enumV})}

\maketitle

\tableofcontents 

\newpage

\part{Lecture and Reading Overviews}

\section{Greek Texts}

\subsection{Greek language basics}

\begin{enumerate}
\item phonology
    \begin{enumerate}
    \item 24 letters: 18 consonants, 6 vowels (see separate sheet)
        \begin{enumerate}
        \item two vowels are always short (epsilon, omicron), two vowels are always long (eta, omega), three vowels can be short or long depending on context (alpha, iota, upsilon)
        \end{enumerate}
    \item syllables can either be long or short
        \begin{enumerate}
        \item long: contains 1) a long vowel or 2) a short vowel followed by two or more consonants
        \item short: contains a short vowel followed by one or zero consonants
        \end{enumerate}
    \end{enumerate}
    \item morphology
        \begin{enumerate}
        \item Greek morphology is notoriously complex, highly inflected
        \item Greek nouns have three genders (masculine, feminine, neuter) and five cases (nominative, genitive, dative, accusative, vocative)
            \begin{enumerate}
            \item a noun's gender is often not recognizable by ending alone
            \end{enumerate}
        \end{enumerate}
    \item meter
        \begin{enumerate}
        \item unlike English meter, which uses stress (emphasis), Greek meter uses quantity, i.e. the length of syllables (see above)
        \item Greek meter consists of a set amount of feet of different forms, with the same definitions as English (switching stress for quantity)
        \item meter was usually rigidly followed according to genre, such that, e.g, epic poetry (Homer) had a different base meter than lyric poetry (Pindar)
        \item meter is often used to determine whether a vowel is long or short
        \end{enumerate}
    \item dialects
        \begin{enumerate}
        \item Greece was a loose collection of independent `city-states', sometimes allied w/ each other, sometimes not
        \item the city-state was the basic unit of allegiance
        \item different areas of the Greek speaking world, at least until somewhat after Alexander the Great, spoke different dialects of Greek, mostly mutually intelligible, but highly important to Greek identity
            \begin{enumerate}
            \item Ionian: Athens, the southwest coast of Turkey
            \item Doric: Sparta, the Peloponnese, northwest Greece
            \item Aeolic: Boeotia, northeast Greece, the northwest coast of Turkey
            \end{enumerate}
        \item different literary genres were associated with different dialects
            \begin{enumerate}
            \item Ionian: epic (Homer, Hesiod), all prose genres, parts of drama
            \item Doric: most lyric, parts of drama
            \item Aeolic: some lyric poetry
            \end{enumerate}
        \end{enumerate}
\end{enumerate}

\subsection{Pindar}

\begin{enumerate}
\item 518-? (at the earliest) 446 BCE, from Thebes
\item Theban context
    \begin{enumerate}
    \item Boeotian dialect, different from Ionic (Athens) or Doric (Sparta)
    \item gives more literary import to Hesiod, rather than Homer
    \item known for religious piety and social conservatism (unlike Athens)
    \item sometimes considered `backwoodsy'
    \item fought with Persia against other Greek city states
    \end{enumerate}
\item from an aristocratic family
\item aristocratic values, unclear politics; seemed to travel around to different games, dependent on wealthy and powerful patrons
\item one of the nine principal/canonical lyrics poets of ancient Greece, often considered the `greatest'
\item epinician odes
    \begin{enumerate}
    \item victory odes for winners of one or some of the four pan-Hellenic games 
        \begin{enumerate}
        \item the Olympian games, held every fours years in of honor Zeus, was the most important pan-Hellenic set of games
        \item winners were highly honored and state-subsidized, always aristocratic
        \item winners would commission the victory odes
        \end{enumerate}
    \item Greek lyric poetry was traditionally sung and danced to, originally religious and communal
    \item Pindar most famous to transpose lyric to athletics but not the first; not a long-lived genre as the games became less important
    \item Pindar has four sets of epinicians organized around each of the four games, ranked by importance of the game; wrote other genres but few survive
    \end{enumerate}
\item style
    \begin{enumerate}
    \item wrote in a literary Doric, rather than his native dialect (common for this genre)
    \item a relatively rolling/free style, compared to Callimachus
    \item a loose and complicated relationship to meter, usually
    \item a loose structure: some kind of intro for the victor and usually in media res, usually uses gnomic (proverbial statements as transitions, some kind of myth retelling, abrupt beginnings/endings
    \item Pindar potentially wrote his own music and choreography
    \end{enumerate}
\item themes
    \begin{enumerate}
    \item communal glory/participation
    \item the nature and benefits of excellence
    \item the goodness of the aristocracy, nobility, grace
    \item the societal function of the poet/poetry
    \item poetic truth and wisdom, often divine
    \item mythic revision in service of piety
    \end{enumerate}
\end{enumerate}

\subsection{Plato}

\begin{enumerate}
\item 427-347 BCE, from Athens
\item Athenian context
    \begin{enumerate}
    \item instituted the first democracy in 510 BCE
    \item rhetoric became extremely important legally and culturally
    \item Peloponnesian War occurred during Plato's lifetime, 431-404 BCE
        \begin{enumerate}
        \item Athens v. Sparta, Sparta narrowly won
        \item Athens former prestige/power greatly diminished
        \item Sparta instituted the 30 tyrants, who were then replaced by a new democratic regime
    \end{enumerate}
    \end{enumerate}
\item often considered the foundation of Western/European philosophy and was an important figure in his own time
\item from an aristocratic family
\item pupil of Socrates
    \begin{enumerate}
    \item famous for his `method', the elenchus, based on rigorous questioning
    \item featured in almost all of Plato's dialogues (Plato himself appears in none)
    \item shifted focus (according to Plato, at least) from metaphysics/cosmology to ethics
    \item executed by the post-war democracy
    \item no extant writings of his own
    \end{enumerate}
\item other influences
    \begin{enumerate}
    \item Heraclitus, known for his metaphysics of constant flux
    \item Parmenides, known for his metaphysics of eternal, unchanging being
    \end{enumerate}
\item writings
    \begin{enumerate}
    \item organized chronologically: early, middle, late
        \begin{enumerate}
        \item \textit{Cratylus} is in the middle period, when Plato starts to break away from Socratic influence/style 
        \end{enumerate}
    \item nearly all in dialogue format, featuring two or more interlocutors in discussion
        \begin{enumerate}
        \item potentially influenced by comedic theater and rhetoric
        \end{enumerate}
    \end{enumerate}
\item themes
    \begin{enumerate}
    \item the Forms, ultimate, unchanging reality that forms the basis of the existing, changing world
    \item eternal soul versus mortal body
    \item recollection via the eternal soul as a means for `learning'
    \item philosophy as a moral imperative
    \item the danger of democracy and the need for aristocratic philosophers
    \end{enumerate}
\end{enumerate}

\subsection{Callimachus}

\begin{enumerate}
\item few certain facts, probably c. 310-240 BCE, from Cyrene
    \begin{enumerate}
    \item Cyrene was an old Greek colony in modern Libya, with a well-known foundation myth
    \end{enumerate}
\item lived/worked in Alexandria
    \begin{enumerate}
    \item Alexandria was founded by Alexander the Great in Egypt, on the coast
    \item becoming a major site of Greek culture and learning around Callimachus' time
    \item part of the Ptolemaic dynasty, one of the successors of Alexander the Great
    \end{enumerate}
\item part of the Hellenistic age
    \begin{enumerate}
    \item literature characterized by small, intricate poems, often opposed to the expansiveness of epic
    \item often meta-reflective on the nature of poetry and the forms that it uses
    \item usually linked to the library of Alexandria and scholarly activity there
    \item active, connected literary community, of which many works survive
    \end{enumerate}
\item works
    \begin{enumerate}
    \item has a large corpus, much of it no longer extant
    \item best known for hymns and short epigrams, which tend to have an extremely difficult and affected style
    \item also a long poem called the \textit{Aitia} (`the causes'), only fragments survive
    \item extensive work in the library of Alexandria, cataloguing works in his \textit{Pinakes} (`tablets')
    \end{enumerate}
\item the hymns
    \begin{enumerate}
    \item a well-established and respected Greek literary form
    \item has a relatively set structure: invocation, prayer, praise
    \item Callimachus consciously imitates this form, two types
        \begin{enumerate}
        \item mimetic (``Hymn to Apollo''): a recreation of the hymn setting/action
        \item non-mimetic (``Hymn to Artemis): narration of the hymn form
        \end{enumerate}
    \item would have been sung or danced to
    \item seem consciously arranged, if not by Callimachus himself
    \item written in an affected Homeric idiom (as was common), despite having a different native dialect (Doric) and the Greek language having changed (Koine)
    \item also develops his own unique vocabulary, frequently uses puns and sonic/visual tricks
    \item poems are highly structured and concise; uses a much more rigid metrical style than Homer and other older poets, usually dactylic hexameter
    \end{enumerate}
\item themes
    \begin{enumerate}
    \item the nature/goals of poetry
    \item the relationship of poetry to royalty and divinity
    \item beauty and its possibility
    \item interpretations of Greek myths
    \item the nature of text, poetry as text, as opposed to oral epic
    \end{enumerate}
\end{enumerate}

\subsection{Sextus Empiricus}

\begin{enumerate}
\item little is known, probably c. 160-210 CE, in a major city of the Roman Greek-speaking world (Athens, Alexandria, maybe Rome)
\item was an important physician and philosopher
    \begin{enumerate}
    \item part of the Empiricist school of medicine, which emphasized experience and observable symptoms
    \item opposed to Rationalist/Dogmatic and Methodic schools
    \end{enumerate}
\item a proponent of and only surviving example of Pyrrhonist skepticism, a primary school of philosophy of the time
    \begin{enumerate}
    \item following from Pyrrho (no writing of his own)
    \item considered `empiric' in a way similar to the medical school
    \item distinct from Academic skepticism, the descendant of Plato's Academy
    \item at odds with the contemporary Stoics and Epicureans
    \item constitutes a procedure for rooting out unjustified belief and suspending judgement/belief for the sake of `tranquility'
    \item i.e., nothing can be decided/believed (rather than, strictly, `known')
    \item distinct from nihilism
    \end{enumerate}
\item works
    \begin{enumerate}
    \item \textit{Outlines of Pyrrhonism}, an intro to skepticism
    \item \textit{Against the Disciplines}, an application of the the skeptical method on various disciplines
        \begin{enumerate}
        \item disciplines conspicuously similar to the medieval curriculum
        \end{enumerate}
    \end{enumerate}
\item themes
    \begin{enumerate}
    \item the finding of irreconcilable differences (impressions, claims) within the same phenomena, unresolvable contradictions
    \item search for / goal of tranquility, undisturbedness
    \item uselessness of explicit/intensive learning, usefulness of prima facie or societal/intuitive knowledge
    \item the development of practical abilities over theoretical learning
    \end{enumerate}
\item literary period was the Second Sophistic, which featured a performative use of older literary/linguistic forms
\end{enumerate}

\newpage

\section{Modernist Texts}

\subsection{Dickinson}

\begin{enumerate}
\item 1830-1886, Amherst, MA
    \begin{enumerate}
    \item prominent, middle-class, puritan family
    \item well-educated
        \begin{enumerate}
        \item unusual for women at the time
        \item prep school, including classics
        \item briefly attended Mount Holyoke Female Seminary
        \end{enumerate}
    \item famously reclusive
        \begin{enumerate}
        \item exclusively in Amherst by 1850s, reclusive by 1860s
        \item talked to people behind her bedroom door by 1867
        \item ``the woman in white''
        \item few friends, avid letter writer: most burned upon her death by her request
        \end{enumerate}
    \item known during her lifetime as a gardener
    \end{enumerate}
\item poems
    \begin{enumerate}
    \item around 1800 poems, only a handful published during her lifetime
    \item first collection in 1890, highly edited
    \item later attempts to restore original version
    \item three ``periods'': pre-1861, 1861-1865, post-1866, most falling in the middle period
    \end{enumerate}
\item style
    \begin{enumerate}
    \item odd, convoluted syntax
    \item idiosyncratic m-dashes and capitalization
    \item ballad stanza
        \begin{enumerate}
        \item quatrains, ABCB rhyme
        \item iambic tetrameter or trimeter
        \end{enumerate}
    \end{enumerate}
\item themes
    \begin{enumerate}
    \item death/mortality
    \item nature/Nature (gardens/flowers)
    \item Deity/Being (transcendentalist influence)
    \item the Master
    \end{enumerate}
\item reception
    \begin{enumerate}
    \item neglected until the 1920s, now a major figure
    \item often contrasted to Walt Whitman: introverted vs. extroverted
    \item considered a ``proto-modernist'', early feminist poet
    \end{enumerate}
\end{enumerate}

\subsection{Mallarm\'{e}}

\begin{enumerate}
\item 1842-1898, Paris
    \begin{enumerate}
    \item worked as a school teacher of English grammar
    \item famous for his salons, gatherings of prominent French writers/thinkers
    \item among the most influential French poets, especially for literary theory
    \end{enumerate}
\item relevant literary movements
    \begin{enumerate}
    \item Romanticism (Hugo): the artist
    \item Realism/Naturalism: novels
    \item Parnassianism (Gauthier): art-for-art
    \item Decadence (Baudelaire): aesthetic depravity
    \item Symbolism (Mallarm\'{e}): aesthetic truth
    \item Modernism (T.S. Eliot): a kind of superset including Decadence through later movements, rooted in ``experimentation''
    \end{enumerate}
\item works
    \begin{enumerate}
    \item relative to his fame, published little
    \item translated Edgar Allen Poe
    \item a couple English grammars
    \item several long poems, sonnets, prose poems, some in collections
    \item most famously, the labyrinthine \textit{Un Coup de d\`{e}s}, \textit{A Throw of the Dice}
    \end{enumerate}
\item style
    \begin{enumerate}
    \item syntactically difficult, formally innovative
    \item canonical symbolism, literary motifs
    \item prone to abstractions 
    \item direct relationship between content and form
    \item emphasis on rhythm / musical qualities
    \item constant word play
    \end{enumerate}
\item themes
    \begin{enumerate}
    \item aesthetics can uncover Truth
    \item the end of poetry
    \item the symphonic nature of poetry
    \end{enumerate}
\end{enumerate}

\subsection{Saussure}

\begin{enumerate}
\item 1857-1913, Geneva
    \begin{enumerate}
    \item from wealthy family of Genevan academics
    \item some early studies in Geneva
    \item 1876-1880, linguistics graduate work at the Universities of Leipzig and of Berlin
    \item 1878, work on Indo-European vowels
    \item taught at University of Paris until 1892, then at the University of Geneva
    \item 1907, begins teaching the Course in General Linguistics
    \end{enumerate}
\item modernizing Europe
    \begin{enumerate}
    \item general: industrialization and conservative movements after the 1848 revolutions
    \item Germany: consolidating under Prussia, making alliances, industrializing
    \item France: second Napoleonic Empire (1852-1870), Franco-Prussian war (1870), Paris Commune (1871), Third Republic (1871-1940)
    \item Switzerland: only recently federalized (1848), strong local allegiances
    \end{enumerate}
\item the field of linguistics
    \begin{enumerate}
    \item relatively new as a scientific field of study (as opposed to the philosophy of language), around the 18th century
    \item focus on comparative linguistics, esp. historical linguistics during the 19th century
    \item Neogrammarians in Leipzig, later 19th century, still focus on diachronic; scientific requirement but narrow scope
    \end{enumerate}
\item works
    \begin{enumerate}
    \item relatively few publications: some technical works on Indo-European
    \item difficulty organizing theoretical insights: major work unpublished, most known from an edited collection of course notes (\textit{Course in General Linguistics}
    \item notes supposedly for a potential book published in the 1990s
    \end{enumerate}
\item key ideas
    \begin{enumerate}
    \item dual character and arbitrariness of the sign
    \item differential nature of the lexicon
    \item linear nature of speech
    \item social structure of language
    \item renewed focus on synchronic, with diachronic: theories of language change/use
    \end{enumerate}
\item influence
    \begin{enumerate}
    \item sets up terms for much of modern linguistics, though he isn't considered scientific
    \item insights incorporated into structuralism and other French literary theories
    \end{enumerate}
\end{enumerate}

\subsection{Stein}

\begin{enumerate}
\item 1874-1946, born in Allegheny, PA
    \begin{enumerate}
    \item upper-middle class Jewish family
        \begin{enumerate}
        \item good/aesthetic relationship with brother, Leo; another brother Michael
        \item mother died when 14, father when 17
        \end{enumerate}
    \item 1877 in Vienna/Paris, 1878 Oakland, CA
    \item 1892, Baltimore, with aunts
    \item 1893-1897, Radcliffe College of Harvard
        \begin{enumerate}
        \item studied under important philosopher/psychologist William James
        \item supposedly did experiments with stream-of-consciousness
        \end{enumerate}
    \item 1897--1901, Johns Hopkins Medical School
        \begin{enumerate}
        \item first lesbian relationships
        \item dropped out from boredom
        \end{enumerate}
    \item 1903-1914, in Paris with Leo, Rue de Fleurus
        \begin{enumerate}
        \item starts collecting art, especially modern, cutting edge works: Renoir, Matisse, Picasso, etc.
        \item becomes known from hosting salons: attended by big expat authors, e.g., Hemingway, Fitzgerald
        \item 1907, meets life partner Alice B Toklas
        \end{enumerate}
    \item 1914-1918, fall out with Leo, WWI
        \begin{enumerate}
        \item partly in London, volunteered as medical/additional aid in French army
        \end{enumerate}
    \item 1918-1938, Paris
        \begin{enumerate}
        \item tries to resume former activities
        \end{enumerate}
    \item 1938-1946, WWII, Bilignin, SE France
        \begin{enumerate}
        \item friends with reactionary French leader
        \end{enumerate}
    \end{enumerate}
\item politics
    \begin{enumerate}
    \item odd mix of progressive and reactionary
    \item not particularly outspoken about explicit politics
    \item openly lesbian, perhaps implicit patriarchal interactions
    \item various racist views on immigration
    \item relatively pro-America, critic of New Deal
    \item helped Allied powers in WWI, friendly with Vichy government in WWII (despite being Jewish)
    \end{enumerate}
\item works
    \begin{enumerate}
    \item many books of various genres, often published later than original compositions
    \item known for recycling previous works, form journals
    \item desired fame, high opinion of herself
    \item 1909, \textit{Three Lives}, first big hit
    \item 1914, \textit{Tender Buttons}, big/canonical ``modernist'' hit
    \item 1933, \textit{Autobiography of Alice B. Toklas}, most popular book
    \end{enumerate}
\item style
    \begin{enumerate}
    \item difficult to parse, paratactic vs hypotactic
    \item highly repetitive
    \item focus on ``ordinary'' language
    \item sparse punctuation
    \item sparse in ``content''
    \item highly edited and re-edited
    \end{enumerate}
\item themes
    \begin{enumerate}
    \item lumped with modernism
    \item the ability of aesthetics to redeem, make anew, solve problems
    \item the need to recreate / revivify words: `` a rose is a rose is a rose is a rose''
    \item difficulty in order to ``shock'', reconfigure expectations
    \item focus on linguistic perceptions
    \item stream-of-consciousness-like sequences
    \end{enumerate}
\end{enumerate}

%\section{Contemporary Texts}

%\subsection{Chomsky}

%\subsection{Plath}

%\subsection{Lorde}

%\subsection{Sun Ra}

\newpage
\part{Notes on Plato's \textit{Cratylus}}

\section{Outline Example}

\begin{small}
\begin{enumerate}
\item Hermogenes initial claim: names are only conventional, 383c-385e
    \begin{enumerate}
    \item horse/man example, 385a
    \item infinite names, including foreign 385d-385e
    \end{enumerate}
\item Socrates contra Hermogenes: names have some natural basis, 385e-386e
    \begin{enumerate}
    \item objects/beings have some natural basis, 385e-386e
        \begin{enumerate}
        \item there are good and bad people, 386a-386e
            \begin{enumerate}
            \item there are wise and foolish people (therefore Protagoras wrong), 386c-386e
            \end{enumerate}
        \end{enumerate}
    \item by extension names have some natural basis, 386e-396d
        \begin{enumerate}
        \item actions are a subset of beings, have a fixed nature, 386e-387b
        \item speech is a  subset of action, 387b-387d
            \begin{enumerate}
            \item truth/falsehood, 387c-387d
                \begin{enumerate}
                \item compositional nature of truth, 387c-387d
                \end{enumerate}
            \end{enumerate}
        \item names are speech and a subset of action, 387d-391a
            \begin{enumerate}
            \item names are tools for instruction (weaving/carpentry examples), 388a-388d
            \item names are made by rule-makers, 388d-390b
                \begin{enumerate}
                \item rule setter uses patterns to embody sounds, 389b-390b
                \end{enumerate}
            \item instructors give specifications to rule-makers, 390b-391a
            \end{enumerate}
        \item judging the correctness of names from Homer, 391a-396d
            \begin{enumerate}
            \item divine versus mortal names (rivers), 391d-392b
            \item divergent human names (Hector, Agamemnon, others), 392b-396d
            \item inheriting names, 393b-395d
                \begin{enumerate}
                \item names can be modified w/o meaning change, 393d-393e
                \end{enumerate}
            \item Zeus etymology 1, 395d-396d
            \end{enumerate}
        \end{enumerate}
    \end{enumerate}
\item digression, Socrates' inspiration, 396d-397a
\item particular etymologies, 397a-427d
\end{enumerate}
\end{small}

\section{Introduction / Hermogenes' Stance, 383a-385e}

\begin{itemize}


\item \textbf{383a}, \textit{Shall we let} --- Socrates starts as an outsider, perhaps a mediator.

\item \textbf{383a}, \textit{Cratylus says} --- Cratylus' argument is set up as the locus of debate, but Hermogenes' side will be tackled first.

\item \textbf{383b}, \textit{He responds sarcastically and makes nothing clear. He pretends to possess some private knowledge ...} --- The idea of naming is immediately set up as a joke, but how coherent Cratylus' joke is isn't clear; i.e., Cratylus is using Hermogenes' name and denying its validity, based on some (supposedly) objective criteria. But Cratylus is also (if Hermogenes is perceiving correctly) not presenting his own position seriously, which might lead to something closer to Hermogenes' position anyway.

\item \textbf{384a}, \textit{oracular utterances} --- The oracular form / divine nature of poetic knowledge is an important theme in Plato, generally.

\item \textbf{384a}, \textit{`fine things' ... names} --- Note that it's not clear or necessary here that names are or be `fine things', nor that Plato himself would make such an equation. The `fine things' is a form of the Greek \ensuregreek{καλός}, which is an extremely important aesthetic and ethical category in Greek and Platonic thought.

\item \textbf{384b}, \textit{fifty-drachma ... one-drachma} --- Even one drachma was a relatively expensive amount for the time. Who gives/gave the one-drachma account is left unspecified.

\item \textbf{384c}, \textit{I don't know} --- An example of a very typical Socratic denial of knowledge.

\item \textbf{384c}, \textit{joint investigation} --- Socrates starts his standard dialogic method for reaching truth.

\item \textbf{384d}, \textit{I believe that any name you give a thing is its correct name. If you change its name and give it another, the new one is as correct as the old.} --- \ul{The initial statement of Cratylus' position}, a strong conventionalist account of names (i.e., the meaning of names are arbitrary).

\item \textbf{385a}, \textit{Whether it is a private individual or a community} --- This is quickly assumed, but it's not even clear that Socrates will end up, at least implicitly, believing this. This distinction will become quite important in Saussure.

\item \textbf{385a}, \textit{horse} --- Horses are a very typical example in Plato and an important symbol of the aristocracy/nobility throughout Western literature.

\item \textbf{385d}, \textit{however many names} --- This might imply that even analytically contradictory names for the same thing are unable to be a contradiction, a perhaps unwelcome implication.

\item \textbf{385d}, \textit{Greeks ... foreigners} --- It's perhaps not insignificant that this example is chosen, given the Greeks' famous xenophobia; i.e., this implicitly states that Greek and non-Greek languages have identical functions or ways of operating, even if their actual words are different.

\end{itemize}

\section{Socrates Refutes Hermogenes, 385e-390e}

\begin{itemize}

\item \textbf{385e}, \textit{whether the same also seems to you to hold of the things that are} --- \ul{Socrates' first lemma to prove}. This is a common move for Plato, to initially base the argument on material objects.

\item \textbf{386c}, \textit{one person to be wise and another foolish ... isn't possible} --- Socrates is doing a proof by contradiction. The idea is that a truly relativistic epistemology (i.e., if everything that exists reduces to what people believe) everyone is equally and maximally wise, so gradations of wise, and therefore foolishness (the least wise), make no sense. This isn't a very solid proof, since it could still be the case that `foolishness' maps onto some set of arbitrary beliefs. 

\item \textbf{386e}, \textit{things have some fixed being or essence of their own} --- \ul{Socrates' first lemma demonstrated}, in accord with Plato's general non-relativistic outlook.

\item \textbf{386e}, \textit{don't actions constitute some one class of the things that are} --- Important link between things and (eventually) words.

\item \textbf{387a}, \textit{to cut in accord with the nature of cutting} --- Plato espouses in his \textit{Phaedrus} the method of `cutting nature at its joints'. Note that, at this point, it would still be an option for any given non-cutting action to be exactly that non-cutting action and therefore `correct' in its own right, a problem that will be explored further later (see 389b).

\item \textbf{387b}, \textit{but with the correct $[$belief$]$} --- Note the sly move from the earlier `accord with the nature of cutting and being cut' to the present `accord with ... the correct \textit{belief}'. This isn't a negation but it implies that belief plays a role in the `nature' of things, a result Plato may not be happy with.

\item \textbf{387b}, \textit{speaking or saying one sort of action} --- The things-actions-sayings progression is complete, meaning that, in some sense, `sayings' have some sort of existence, via extension or subset of actions more generally.

\item \textbf{387c1}, \textit{the natural way to say them ... accomplish nothing} --- The strongest implication here, though not explicit, is that, without some outside-of-word criteria, language would be pure gibberish. The view here, moreover, doesn't necessarily entail that there is strictly only one `natural' way to speak.

\item \textbf{385b2}, \textit{some statements are true, while others are false} --- A foundational belief for Plato and probably for epistemology in some abstract sense.

\item \textbf{385b2}, \textit{those that say ... false} --- The action of `saying' is looped back around to things that exist, which would already contradict Hermogenes' conventionalist position.

\item \textbf{385c}, \textit{it is possible ... in a statement} --- A strict criteria of truth based on compositionality, namely that for a whole to be true all of its parts need to be true. This criteria is not usually assumed in modern logic and is in fact assumed to be false for standard conditionals (i.e., in modern logic, [A $\rightarrow$ B] $\nrightarrow$ [A $\wedge$ B])

\item \textbf{385c}, \textit{smaller than a name} --- The status of letters or phonemes remains unclear so far.

\item \textbf{385c}, \textit{and a part of a false statement is false} --- This is a surprising and even more strict definition of the compositionality of truth. One might expect that a whole could be false if only (at least) one of its parts were false, even if the others are true, but this is not the case.

\item \textbf{385c}, \textit{possible to say ... possible} --- Names can be said separately from statements and have truth value therein.

\item \textbf{387c6}, \textit{we must name them in the natural way for them to be named} --- \ul{Hermogenes' position is essentially refuted}, insofar as a `natural' basis for names is conceded.

\item \textbf{387e}, \textit{weaving} --- A particularly important metaphor for poetic production throughout Greek literature, including Homer and Plato.

\item \textbf{387e-388a}, \textit{what one has to name ... a name is also a sort of tool} --- \ul{The beginning of Socrates' own positive argument as to the nature of names.} An important, three-fold distinction is implicit here, namely that the nature of the thing cut/named (wood/being), the action of cutting/naming, and the tool for cutting/naming (a bit/name) all have three, distinct natures. The middle term, i.e., the action, is perhaps the most surprising and is not at all intuitive, even for the non-name examples, but it was essentially assumed to be the case (above). At least with cutting wood, though, it's easy to imagine a drill that never `partakes' in cutting. On the other hand, a similarly inoperative name is extremely unclear, which is why it would be easy to elide the act of naming and the name itself (as in Hermogenes' first position). I.e., it's important to remember that these are distinct things in Plato's system, and that named-naming-name have separate truth values.

\item \textbf{388b}, \textit{instruct each other ... divide things} --- See the note on 387a above. This equation sets up Socrates'/Plato's dialectic method as the proper naming activity. It remains an open question whether there are non-`dividing' types of activity.

\item \textbf{388c}, \textit{instructor} --- This instructor figure, who uses (pre-built) names, is somewhat slyly inserted here.

\item \textbf{388c}, \textit{those who possess the craft} --- That certain skills/traits belong to certain people/things is a crucial point for Plato, especially in \textit{The Republic}.

\item \textbf{388d}, \textit{rules provide us with $[$names$]$} --- The Greek word for `law' here is \ensuregreek{νόμος}, one of the words Hermogenes used for `convention'. It's surprising that Socrates makes this claim, i.e., that `rules' in some way play a role, even though he already refuted the idea that `convention' is not the criteria for naming.

\item \textbf{388e}, \textit{the one who possesses the craft} --- At least, unlike Hermogenes' first position, only certain people can, by convention, establish names.

\item \textbf{389a}, \textit{that sort of thing whose nature is to weave} --- Note that the shuttle-maker doesn't look to the things to be woven (i.e., thread), but to (the form of) the shuttle itself. A distinction is made here, then, between instances of a shuttle/name and the form/pattern for those instances, another common Platonic move.

\item \textbf{389b}, \textit{suppose the shuttle breaks} --- Particular instances of forms can be faulty, an important point insofar as it confirms the not-strictly-conventional nature of Socrates' position. A strictly conventional position would not allow things/names to be faulty in any sense, since they would precisely describe any one thing, itself.

\item \textbf{389b-c}, \textit{what a shuttle itself is ... how to embody} --- The difference between the form itself and its instances is solidified.

\item \textbf{389d}, \textit{embody in sounds and syllables ... same syllables} --- \ul{An important clarification for the compositional nature of truth.} While names are the most basic units of truth (385c), letters are the material for those units. Letters, apparently like iron as a material, cannot be true or false in the same way as words or shuttles. This is necessary, if unclear, for Socrates' position here, in line with his compositional definition of a statement's truth. Otherwise, only one word, across all languages, could be the correct set of syllables for any given thing it names. A `name itself' then is some generalized form, and one not necessarily unique for each name.

\item \textbf{390a}, \textit{Greek and foreign} --- Again, the equality of Greek and non-Greek languages is somewhat unorthodox for the time.

\item \textbf{390b}, \textit{ship-builder .. captain} Ship building/steering is another consistent metaphor in Greek thought and in Plato, and it carries overtones of hundreds of years of going Greek colonization across the Mediterranean.

\item \textbf{390c}, \textit{supervise the work of a rule-setter} --- Another instance where the analogy is intuitive w/r/t carpenters but not the rule-maker. We haven't been give a true sense of who or what a rule-maker might be, but apparently the rule-user (the namer, the dialectician) can interact with them/it. Furthermore, a carpenter can be both a maker (of shuttles) and a user (of drills), but it's not clear that the rule-maker can be a user or that the rule-user can be a maker (unless the roles refer to the same entity).

\item \textbf{390c}, \textit{ask questions ... dialectician} --- \ul{Socrates' method is explicitly made the best form for understanding names.} Is asking questions tantamount to making statements (and therefore also instances of using names)?

\item \textbf{390d-e}, \textit{Cratylus is right ... isn't easy for me} --- \ul{Socrates affirms Cratylus' position (for now)}, but Hermogenes remains unconvinced.

\end{itemize}

\section{Etymological Examples, 390e-428d}

\begin{itemize}

\item \ul{\textbf{all} --- The etymologies in this section have varying degrees of sensibility. Some of them are basically nonsense, though Plato is likely in on the joke.} See \href{https://plato.stanford.edu/entries/plato-cratylus/#Ety}{\ul{this article}} for a summary of the etymological principles invoked. Almost all of the etymologies take a Heraclitean perspective as their basis, i.e., that the universe is in constant flux, a position Plato strongly disagrees with.

\end{itemize}

\subsection{Etymology 1: Human Names from the Poets, 390e-396e}

\begin{itemize}

\item \textbf{391a}, \textit{showed me ... consists in} --- Hermogenes resists Socrates' conclusion and asks for examples.

\item \textbf{391a}, \textit{I don't have a position} --- Another Socratic denial of knowledge, reaffirming the supposedly joint (dialectical) nature of the question/method at hand.

\item \textbf{391b}, \textit{our next task} --- \ul{Namely, to figure out what a correct name is.}

\item \textbf{391c}, \textit{the sophists} --- The sophists are Plato's canonical rivals, and Plato devotes several dialogues to undermining their very real societal and theoretical authority. The joke here seems to be that sophists are conventionally given the status of being able to speak well, i.e., to name well, and Protagoras, as a relativist, would be especially adept at naming. This relativism has already been rejected, though, as Hermogenes notes. Thus, another rejection of standard convention is enacted insofar as those who are conventionally thought of as good namers are also rejected.

\item \textbf{391c}, \textit{Homer and the other poets} --- Plato famously pits philosophy against poetry, making this somewhat dangerous territory for Plato. Still, the poets are the more traditional namers and the sophists the newer namers. Homer was essentially the basis for Greek education and his importance for Greek society at the time couldn't be overstated.

\item \textbf{391e-392b}, \textit{the gods call things ... more within human power} --- Standard Socratic piety w/r/t the infallible nature of the gods and restriction of the Socratic method to human (ethical/epistemic) affairs.

\item \textbf{392c}, \textit{wiser ... more foolish} --- Good/bad were mapped onto wise/foolish (386b), so there is an implicit connection between those who are good/bad and those who name well/poorly.

\item \textbf{392c}, \textit{a city's women or its men} --- Greek society was famously patriarchal. That the `city', the basis for Greek political life, is introduced here is interesting.

\item \textbf{392d}, \textit{more correct} --- Just like there can be poorly made names, the correctness of names is not necessarily binary.

\item \textbf{393a}, \textit{a man possess that of which he is lord ... I don't think you're wrong} --- A perhaps tongue-in-cheek explanation that is almost certainly based on faulty logic from Socrates (that a lord `possesses' doesn't imply that a `possessor' is a lord), but somehow it gets reaffirmed by Hermogenes (thus, the danger of dialectics).

\item \textbf{393c}, \textit{contrary to nature} --- A thing's unique/observed characteristics are more important than its shared/expected characteristics.

\item \textbf{393c}, \textit{trick ... same argument} --- Socrates does, in fact, seem to be tricking Hermogenes. The argument is based on the fact that kingship is naturally hereditary, so a non-natural offspring of a king would be a non-king, which is exactly what Astyanax, Hector's son, turns out to be, since the city of Troy is destroyed before Astyanax can become king.

\item \textbf{393d-e}, \textit{same thing is signified ... expressed in its name} --- \ul{A reiteration of the relative arbitrariness of letters, distinct from the being/essence/truth of names.}

\item \textbf{393e}, \textit{as long as we include the force or power} --- Again, the truth of names is not dependent on every part of a name (each letter) being true, unlike the relation between names and statements.

\item \textbf{394a}, \textit{seem different to the uninitiated} --- The underlying unity to some superficial discrepancy is a classic Platonic stance. Note also that a correct name doesn't necessarily imply that the name accounts for all of a person's attributes.

\item \textbf{394c}, \textit{same force or power} --- The stable entity is now `power', as opposed to forms or natures.

\item \textbf{394c-d}, \textit{shouldn't have his father's name} --- It's unclear whether this means the natural power should be known at birth or that the person should be renamed upon learning his true power.

\item \textbf{395a-c}, \textit{Agamemnon ... Pelops} --- In these examples, the attributes signified by the names are strictly not inherited.

\item \textbf{395e}, \textit{disguised the name} --- This calls into question how serious Socrates is being here, since it would seem to open up the possibility that names can be reverse engineered as correct after assuming a series of `disguise' maneuvers.

\item \textbf{396b}, \textit{offspring of a great intellect} --- Zeus was often considered to be all-knowing.

\item \textbf{396c}, \textit{astronomers ... great intellect} --- Probably a joke on the `first' philosopher Thales, who was said to have fallen into a well while looking up at the sky.

\item \textbf{396d}, \textit{I do not know ... like a prophet} --- Another common form of the Socratic denial of agency.

\item \textbf{396d}, \textit{Euthyphro} --- Plato has another dialogue called \textit{Euthyphro}, in which Socrates talks to the titular character, who was a priest, just before himself going to trial in Athenian court. Socrates is tried and sentenced to death for `corrupting the youth' and `atheism'.

\item \textbf{396e}, \textit{tomorrow ... purify} --- Socrates again questions the validity of this process.

\end{itemize}

\subsection{Etymology 2: Particular Names, 396e-410e }

\begin{itemize}

\item \textbf{397b}, \textit{things that by nature always are} --- A typical statement of Platonic interest in unchanging things, an assumption/goal perhaps at odds with the discussion to follow.

\item \textbf{397d}, \textit{many foreigners} --- A stereotype held by Greeks about some non-Greeks, especially Persians, e.g., in Herodotus.

\item \textbf{398a}, \textit{he too belonged to the golden race} --- That a contemporary human could be part of the `golden' race makes no sense according to Hesiod's genealogical account, so this is an instance of Socrates rationalizing the poets, as he is wont to do.

\item \textbf{398b}, \textit{Old Attic} --- Athenians considered themselves to be autochthonous, born directly from Attic soil, and the Attic/Ionian dialect had special status by virtue of being the dialect of Homer and Hesiod.

\item \textbf{399a}, \textit{too wise} --- Cleverness is seen by Plato as a trick of the sophists.

\item \textbf{399b}, \textit{a phrase ... became a name} --- This is actually a well-established linguistic phenomena, though not for this example

\item \textbf{399c}, \textit{observes it closely and reasons} --- Plato makes a bid for rationality being the fundamental human property.

\item \textbf{399c}, \textit{take the reigns} --- Hermogenes starts to take control of what he wants Socrates to analyze, which is somewhat unusual among Platonic dialogues.

\item \textbf{401a}, \textit{there isn't anything} --- the `self-sustaining' soul is another Platonic trope, e.g., in \textit{Phaedrus}.

\item \textbf{400b}, \textit{it sounds funny} --- More evidence that Socrates might not be taking this seriously.

\item \textbf{400d-401a}, \textit{they call themselves ... know no others} --- A statement about human limits. Humans have no access to divine language and are unable to learn it.

\item \textbf{401b}, \textit{Hestia} --- The goddess of the hearth, often used as an actual or metaphoric center of the earth/universe and made explicit later as the `essence of things'.

\item \textbf{401c}, \textit{lofty thinkers and subtle reasoners} --- Another dig at Thales and other similar thinkers, including Socrates as portrayed by the comic playwright Aristophanes.

\item \textbf{401d}, \textit{Heraclitus} --- An important pre-Socratic philosopher who plays an important foil role in this dialogue, so far only implicitly and to a much greater degree later on.

\item \textbf{401e}, \textit{swarm} --- Bees play an interesting role in poetry, especially as bearers of nectar for Demeter. (Cf. Callimachus' ``Hymn to Apollo''.)

\item \textbf{402a}, \textit{Heraclitus ... Homer} --- Implicit equation between poetry and the philosophy practiced by Heraclitus, both usually rejected by Plato.

\item \textbf{401a}, \textit{Kronos} --- It's unclear if this coexists or replaces the previous etymology for `Kronos'.

\item \textbf{402e}, \textit{Poseidon} --- God of the sea and of horses. Note how Socrates doesn't commit to any of these possibilities.

\item \textbf{403a}, \textit{Pluto} ---God of the underworld/wealth. In Plato's \textit{Phaedo}, Socrates claims that money is the cause of all war.

\item \textbf{403c}, \textit{desire is far stronger} --- Desire is an important theme/category in Plato, especially in his \textit{Phaedrus} and \textit{Symposium}.

\item \textbf{403d-404a}, \textit{desire to associate ... better man ... agitation} --- The validity of this desire/agitation is one of the principle subjects of Plato's \textit{Phaedrus}.

\item \textbf{403e}, \textit{Sirens} --- Sirens sing (i.e., make names) so beautifully and provide such knowledge that they attract and shipwreck sailors.

\item \textbf{404d}, \textit{people nowadays attach ... truth} --- Perhaps referring to people's preference for poetry/sophists over philosophy.

\item \textbf{405a}, \textit{four powers} --- Cf. Callimachus' ``Hymn to Apollo''.

\item \textbf{406a}, \textit{a single-minded ... together} --- A kind of ridiculous tour-de-force of this etymological method.

\item \textbf{406c}, \textit{the playful} --- Socrates becoming explicit about how seriously he takes this etymological method.

\item \textbf{407a-b}, \textit{Athena} --- Goddess of war and wisdom, an extremely important/controlling character in Homer's \textit{Odyssey}. 

\item \textbf{407d}, \textit{horses of Euthyphro} --- See the book's note on this.

\item \textbf{407d}, \textit{Cratylus says} --- A call back to the very beginning, when Cratylus makes fun of Hermogenes.

\item \textbf{407e}, \textit{Hermes ... I'm no Hermogenes} --- The messenger god who helps Odysseus in Homer, known for deception and trickery. Note the meta-textual aspects of this final god example. If Hermogenes is correctly named he is a deceiver, which would imply that he is a deceiver about his name and is thereby incorrectly named. On the other hand, if Hermogenes is incorrectly named, he is a deceiver and is thereby correctly named. Hermogenes goes so far as to disavow his own name. This is, probably, the essential paradox that sits at the bottom of this entire etymological method.

\item \textbf{408c}, \textit{tragic ... falsehoods} --- A perhaps Hesiodic pessimism.

\item \textbf{409a}, \textit{recent theory about the moon} --- Etymological reasoning is implicitly made equal to, if not better than, the naturalistic (as distinct from Socratic) reasoning by those like Anaxagoras. Socrates in Plato's \textit{Phaedo} claims to be disappointed with a book by Anaxagoras.

\end{itemize}

\subsection{Etymology 3: Derived Philosophical Names, 411a-420e}

\begin{itemize}

\item \textbf{411a}, \textit{no inconsequential class of names} --- The following are important words for a Platonic conception of ethics and epistemology.

\item \textbf{411b-c}, \textit{Most of our wise men ... coming into being} --- An attack, via subjectivity and lack of rationality on the Heraclitean view of the world. The etymological method continues under these assumptions, though, an assumption which allows these particular derivations in the first place.

\item \textbf{411d}, \textit{longing for the new} --- In political terms, the `new' is often in Greek thought (negatively) associated with change/disruption in/of political order.

\item \textbf{413a}, \textit{I persisted in gently asking ... the just} --- Socrates is often (self-avowedly) a gadfly figure who won't let common (mis)conceptions remain unexamined. This is a hint at how Socrates will want to account for names separate from these spurious etymologies.

\item \textbf{411b-d}, \textit{the just is the sun ... reasons we mentioned} --- This narrative is also found in Plato's \textit{Phaedo}, of Socrates being dissatisfied with prior/contemporary naturalistic explanations.

\item \textbf{414c}, \textit{name-giver has imitated} --- Imitation is an important category for Plato, acting as the relationship between varying degrees of reality (including, perhaps, the form of a name and the name itself); its status here, at this point, is not clear but will become important later. 

\item \textbf{414b-c}, \textit{if you remove ... sticky ramp} --- A particularly ridiculous etymology, which Hermogenes notes.

\item \textbf{414c}, \textit{dress them up} --- The idea of ornamental features, as opposed to fundamental features, is important for Plato.

\item \textbf{414c}, \textit{time} --- Linguistic theory, generally, must deal with the function of time w/r/t human affairs. Saussure especially emphasizes this.

\item \textbf{414c-d}, \textit{the work of people ... mouths make ... no human being can understand} --- Does this mean the eventual names are untrue? Socrates waffles between ascribing and not ascribing intention to the sound changes that names undergo. See 387c, in which the improper namer `accomplishes nothing'.

\item \textbf{414d}, \textit{add whatever we like} --- See the note on 395e.

\item \textbf{414e}, \textit{supervisor, like yourself} --- Irony on the part of Socrates, since Hermogenes is basically a useless yes-man.

\item \textbf{415c}, \textit{courage ... shackles} --- Socratic thought places a strong emphasis on not fearing death, in accordance with the soul's immortal nature.

\item \textbf{415b-d}, \textit{I don't yet understand ... the name} arete --- Socrates, supposedly, derives in realtime the etymology for the word for `courage', showing either the validity or non-validity of the method. The etymology is based on the derivation of the word for `evil', which is patently nonsense, insofar the derivation consists of taking that same word and adding another, completely different word, to it and then appealing to a foreign language, Socrates' `device'.

\item \textbf{416b-d}, \textit{what caused ... welcome as such} --- The meaning of this derivation seems unclear. It would seem that Socrates stumbles on a derivation for `fine things' from the act of `naming' itself, but he then introduces two other claims: that `thought' `names' and that `naming' is `good' in some way. What this does, though, is set up the name-giver (i.e., the Socratic dialectician (388b)) as someone who does `fine' deeds.

\item \textbf{417b}, \textit{retailers} --- Socrates was known to hang out in the agora, the marketplace, with merchants, as in the \textit{Apology}.

\item \textbf{418a}, \textit{I'm not responsible} --- Socrates reminds Hermogenes how little invested in this work he is.

\item \textbf{418b}, \textit{Fine modern language} --- Perhaps referring to the sophists, or to sympathizers of democracy (which Plato argues against elsewhere).

\item \textbf{418c}, \textit{especially the women} --- Apart from the general Greek stance that language is a `manly' domain, why Plato ascribes this particular quality to women is unclear and is, it turns out, empirically false, insofar as women, relative to men, are usually on the vanguard of linguistic innovation.

\item \textbf{420b}, \textit{opinion ... how things are} --- This first derivation for `opinion' is more likely the opposite of how Plato would describe the yearnings of the soul, since the soul yearns for `truth' instead of mere `opinion'.

\item \textbf{420c}, \textit{thinking is the motion} --- This is perhaps meant to represent another view of Anaximander, as well reference a Greek pun about life/bows made by Heraclitus. Again, the status of `thought' here is not exactly clear.

\item \textbf{420e}, \textit{strength lasts} --- Etymological work is described as something that takes energy. The relationship between the Socratic method/activity to leisure, in general, ambiguous.

\end{itemize}

\subsection{Etymology 4: Primary Philosophical Names, 421a-427d}

\begin{itemize}

\item \textbf{421a}, \textit{name} --- Even the name of `name' is derived and must be fitting to its content. Curiously, the derivation seems to indicate, like the other words, that, as the discussion shows, it takes effort to find names. 

\item \textbf{421c}, \textit{hammered} --- Note the carpentry/blacksmith metaphor, likening statements, again, to tools.

\item \textbf{421d}, \textit{the basic or `first' names} --- A qualitatively different set of names from all the foregoing is being referenced here, the primary ones that can't be broken down into their parts.

\item \textbf{421d}, \textit{plausible ... investigate} --- Plato elsewhere rejects the criteria of the plausible, as here Socrates doesn't rest on that criteria.

\item \textbf{421e}, \textit{give up} --- For Plato, there are final truths, even if humans must continually analyze further and can't reach them.

\item \textbf{422b-c}, \textit{different manner ... correctness} --- Just like primary names compared to derived names are different in structure, so their analysis must be different. The criteria of correctness is the same, however, for both types.

\item \textbf{423a}, \textit{our body imitates} --- The important concept of imitation is brought up again.

\item \textbf{423c}, \textit{imitates something with his voice ... to say at all} --- The value of letters is finally explored, namely that they are imitations of things or ideas. However, this imitation does not constitute naming itself, which means that names can remain the basic constituents of truth value proper.

\item \textbf{423d}, \textit{music and painting} --- The composition of primary names is and will be compared to the other arts.

\item \textbf{424d}, \textit{well-divided} --- Recall that the main function of the name-giver is to divide essential natures (388b).

\item \textbf{424d}, \textit{using one letter for one thing} --- The compositionality of letters is, again, very different from the compositionality of names.

\item \textbf{425a}, \textit{important, beautiful, and whole} --- The criteria for a fully formed word now appears quite different from truth value as Socrates earlier maintained, but presumably the latter still obtains. Namely, a `picture' as a whole, as a set of various imitations, still must have some claim to truth or falsehood. This is a perhaps unintuitive take on aesthetics for modern ears.

\item \textbf{425d-426b}, \textit{deus ex machina ... will be worthless} --- The \textit{deus ex machina} was an actual machine used for plot-saving divine entities in ancient theater. Socrates rejects here, as often, handwaving techniques, especially those that make use of the gods in undefined ways. Socrates also here, as often, espouses a strict criteria of knowledge, namely that any knowledge of a phenomenon requires complete knowledge of it.

\item \textbf{427c}, \textit{make a sign ... things they name} --- In order for this analogy to color in paintings to be apt, it would have to be the case that different letters change by degrees (such that, e.g., slightly different pronunciations of the `same' letter would be significant) and thereby be turned into other letters across a continuum. Plato doesn't explore this issue, but it is, at least for modern linguistics, empirically false.

\item \textbf{426c-427d} --- The idea that certain sounds of relatively fixed connotations is a common one, even today. It's not generally thought to obtain in contemporary linguistic theory, though some relatively recent theorists maintain it, e.g., Jakobson.

\end{itemize}


\section{Cratylus' Stance, 427d-435d}

\begin{itemize}

\item \textbf{427d}, \textit{Cratylus disagrees} --- Socrates asks Cratylus for his opinion because, presumably, Hermogenes could listen to etymologies forever, but this is not the sort philosophical discussion that Socrates is interested in.

\item \textbf{427e}, \textit{learned so quickly} --- Cratylus makes a jab at the etymologies. Socrates supposedly was inspired to give them on the spot, but he must have taken a long time to learn them.

\item \textbf{428c}, \textit{after my own mind} --- Socrates has to force Cratylus to explain his position

\item \textbf{428d}, \textit{self-deception} --- A fundamental category for Platonic thought, though perhaps not a clear one. At points in Plato's work the position that someone can't knowingly do wrong is espoused.

\item \textbf{428d}, \textit{backwards and forwards} --- A basic statement of the Socratic/Platonic method, namely the toggling between levels of phenomena. Found in, among others, the \textit{Phaedo} and \textit{Republic}.

\item \textbf{428e}, \textit{names are spoken} --- It's not clear whether Plato wants to make a distinction here between the giving of a name and the speaking of a name. Rule-setters from above referred to those who make the name, however.

\item \textbf{429a}, \textit{painters} --- Note the shift here from the earlier discussion. Rule-setters were at first likened to weavers and carpenters, now to painters. This difference is less one of kind than our ears might think (since all were `techne', arts), but it does indicate an important distinction with respect to levels of phenomena in Platonic thought, as in the \textit{Republic}.

\item \textbf{429b}, \textit{some names have been better} --- This is a further departure from the original discussion of shuttles on the part of Cratylus, which implies that names can be faulty (398b).

\item \textbf{429b-c}, \textit{rule-setters ... better given} --- Note that rule-setters (those who make names) is slyly elided into those who `give' names, a much more unclear status. This ambiguity persists through the rest of the argument.

\item \textbf{429c}, \textit{hasn't been given it at all} --- \ul{The basic statement of Cratylus' position}, an extremely unintuitive position to take.

\item \textbf{429c}, \textit{is it even possible} --- Socrates is trying to expose the same problem for Cratylus as he did for Hermogenes (387c1), namely that this extreme position implies that nothing can be said (contrary to the bare fact of the discussion at hand).

\item \textbf{429d}, \textit{not saying things that are} --- This might be seen as an extreme implication of the Platonic theory of forms (that every quality has a transcendent form) or of the earlier Socratic ethical position above (that no on can knowingly act falsely), which Plato, by putting it into Cratylus' mouth, will try to defend against.

\item \textbf{430a}, \textit{acting pointlessly} --- Again, an unintuitive response that is easily shown to be false, only by pointing out that, e.g., Cratylus could and likely would respond, if with confusion. This is the critique that Socrates eventually gets to, after some subarguments

\item \textbf{430a}, \textit{a name ... thing it names} --- This distinction is maintained from the earlier discussion.

\item \textbf{430d}, \textit{assignment of a painting or a name} --- The analogy seems to get even further muddled here. Previously the construction of names was like that of a painting, meaning painters are like name-makers. Now whole paintings and whole names seem to be being assigned, which refers not to the making of names but the assigning of names, two tasks that Socrates was earlier at pains to distinguish. It might be the case that Socrates is saying that `likeness' of men or women in painting, and therefore in names, can refer to things like colors. E.g., a painter could paint a `man' despite giving him the (supposed) skin color of a `woman'. This option would recall the natural-non-natural offspring issue from above (393c).

\item \textbf{430e}, \textit{bring before the sense} --- Note the definition of humans from 399c. 

\item \textbf{431b}, \textit{combination of names and verbs} --- This definition of propositional content is still the basic, accepted one.

\item \textbf{431c}, \textit{produces a bad one} --- Gradation in the efficacy of names is brought back.

\item \textbf{431e}, \textit{rule-setter} --- Referring to the person who makes names, which later (435d) becomes confused.

\item \textbf{431e}, \textit{some rule-setters} --- \ul{The basic linguistic theory here is that names aren't completely arbitrary, insofar as not everyone can create them, but they also aren't completely fixed, insofar as a subset of humans do create them.} This is in accord with Plato's supposed political views that espouse rule by a select elite (the aristocracy, philosophers). 

\item \textbf{431e}, \textit{add, subtract ... don't write it at all} --- This is a direct critique of the long etymological discussion Socrates just performed. It seems, also, though, to be a reversal of Cratylus' position to something more like Hermogenes' position.

\item \textbf{432b}, \textit{an image cannot remain an image} --- A fundamental Platonic position, that images are qualitatively different from what they depict. Specifically, images are unable to be perfect, but always fall short.

\item \textbf{432b}, \textit{some god ...painters do} --- The gods frequently made themselves into the likeness of various, real or fake, humans in Greek poetry and myth. Socrates likens the gods to painters.

\item \textbf{432d}, \textit{no on would be able to say} --- Again, the danger is that names are completely useless.

\item \textbf{432e}, \textit{things are still named} --- \ul{Socrates has refuted Cratylus' most extreme position}: names are not completely useless, even if not perfect (Hermogenes would turn around if called).

\item \textbf{433a}, \textit{if it includes its pattern} --- Recall the discussion of shuttles from above. The 'pattern' is an important term for Plato but not same word translated as `pattern' here.

\item \textbf{433b}, \textit{by means of letters and syllables} --- \ul{A confirmation of the compositional nature of names}.

\item \textbf{433d}, \textit{I'm not satisfied} --- Cratylus still isn't convinced, and Socrates ends up making essentially the same argument as before except he reroutes it through Hermogenes' prior position.

\item \textbf{433e}, \textit{already knew the things} --- This makes more explicit the importance of name-giving. Not everyone knows the nature of things, and names are useful for imparting that knowledge, hence Plato's use of `instructors'.

\item \textbf{434d}, \textit{the opposite of hardness} --- This also acts somewhat as a critique of the etymological method from before.

\item \textbf{435a-d}, \textit{if that's right ... opposite kind of names} --- Convention still plays a role in Socrates' position (perhaps one larger then he wants), but not indefinitely like Hermogenes' prior position. Usage seems to be a continued use of some given convention. Speakers know what each other `mean' when they speak, even when they speak incorrectly somehow. This is, actually, a very difficult empirical fact to explain.

\end{itemize}

\section{Names and Knowledge, 435d-439b}

\begin{itemize}

\item \textbf{435d}, \textit{what power} --- Socrates assumes his point is proven and moves onto what he wants to prove. However, a confusion happens here. So far, with Cratylus, Socrates has been talking about the name-maker, the rule-setter, but now he switches to `instruction', which is the rule of the namer or name-user, a distinct entity. \ul{These two formerly distinct roles get conflated into one role: the name-giver.}

\item \textbf{435d-e}, \textit{when you know ... same craft} --- Note that the syllogism might not be too helpful. Even if it holds that [1) knowing what a name is like] and [2) the name is like the thing it names] implies [3) one also knows the thing], it does not follow that [4) one knows that the name is like the thing it names]. However, the `one craft' note doesn't mean that the name and named thing are the same, just like wood and drills are different but fall under carpentry.

\item \textbf{436a}, \textit{discover the things} --- I.e., Socrates'/Plato's dialectic is the best way to discover the `things that are', i.e., the forms, the epistemological ground of thought.

\item \textbf{436b}, \textit{being deceived} --- Typical Socratic distrust of tradition and of how things are normally thought.

\item \textbf{436c}, \textit{otherwise ... names at all} --- Cratylus refuses to back down from this position

\item \textbf{436c}, \textit{consistent with one another} --- To use modern terms, precision in no way entails accuracy, as Socrates notes.

\item \textbf{436d}, \textit{every man must} --- For Socrates/Plato, philosophy was an ethical imperative, even if not everyone could perform it.

\item \textbf{436e}, \textit{all things are moving} --- \ul{Socrates now can start attacking the Heraclitean assumptions directly.}

\item \textbf{437c}, \textit{very worst ... very best} --- Note that this doesn't support what would be a Platonic view against the Heraclitean view. Instead it shows that the method itself is flawed.

\item \textbf{437d}, \textit{count names like votes} --- Playing on assumed anti-democratic leanings amongst the elite.

\item \textbf{438b}, \textit{impossible to learn ... ourselves} --- An important disjunction: names can either be learned from others or given directly. If things can only be known through their names (as Socrates assumes here ad hoc, per `on our view'), this disjunction entails a contradiction, since there would no way for the first names to arise. \ul{This allows Socrates to claim that knowledge apart from names is possible}, inkeeping with the general Platonic position that affirms the existence of non-linguistic entities.

\item \textbf{438c}, \textit{contradicted himself} --- That the divine beings can't be or make contradictions in any way is exploited in the \textit{Republic}.

\item \textbf{439e}, \textit{independently of names} --- What constitutes this non-name learning is unclear and will remain so.

\item \textbf{439e}, \textit{through one another ... themselves} --- The benefits of two like minds is a consistent trope in Plato and is essentially the basis of the Socratic dialectic. Still, what learning `through oneself' constitutes is unclear, though it would presumably be what Plato, in some sense, does.

\item \textbf{439b}, \textit{better to learn} --- Another classic Platonic position that requires searching for truth `itself'.

\end{itemize}

\section{Forms and Conclusion, 439b-440e}

\begin{itemize}

\item \textbf{439c}, \textit{aren't really that} --- Socrates takes on Heraclitus directly.

\item \textbf{439c}, \textit{dream about} --- Socrates often appeals, especially near the ends of dialogues, to dream or other irreal states for his larger, metaphysical positions.

\item \textbf{439c}, \textit{beautiful itself} --- A basic statement of the Platonic theory of forms, the transcendent entities.

\item \textbf{439e}, \textit{be something at all} --- The most basic contradiction / problem that Plato finds for the in-flux philosophies. The proof here is not very strong, since it starts out assuming what it wants to show, namely, that `beauty' is always the same.

\item \textbf{439e}, \textit{can't even be known} --- A much stronger proof of the above, rooted in an epistemological contradiction.

\item \textbf{440a}, \textit{no such thing as knowledge} --- \ul{Plato's basic position, that philosophical inquiry is impossible, i.e., ends in contradiction, if based on some conception of flux.}

\item \textbf{440c-d}, \textit{all flowing ... drip over} --- A set of metaphors making fun of Heraclitean positions by making them trivial/gross instead of regal/philosophical.

\item \textbf{440d}, \textit{still young ... more as Heraclitus} --- Perhaps an explicit `teaching' moment, which Cratylus rejects.

\item \textbf{440e}, \textit{into the country} --- Perhaps meant to signify that Cratylus is going into a pastoral setting, a literary trope that features an idealized, unchanging, and untainted rustic landscape.

\item \textbf{440e}, \textit{Hermogenes will} --- Hermogenes, ironically, seems to take on his namesake afterall.

\end{itemize}

\newpage
\part{Notes on Sextus Empiricus' ``Against the Grammarians''}

\section{Introduction to the Entire Work}

\begin{itemize}

\item $[1]$, \textit{the counter-argument} --- Note that this book, which is a largely a series of counter-arguments, starts off with a counter-argument against a counter-argument, the latter of which comes from a different, rival philosophical sect

\item $[2]$, \textit{Pyrrho's pupil}, Epicurus was a disciple of Nausiphanes who in turn was a disciple of Pyrrho. Epicurus, according to Sextus, repudiates his Pyrrhonic lineage; however, Sextus himself is a `Pyrrhonian skeptic'

\item $[3]$, \textit{self-taught philosopher} --- Perhaps a reference to the Platonic idea that the true philosopher interacts with / learns from the eternal Forms.

\item $[5]$, \textit{dogmatic ... mob} --- The skeptic, purportedly tries both to avoid prescriptive, study-based claims, as well as democratic, persuasion-based claims.

\item $[6]$, \textit{lack of uniformity} --- Perhaps Heraclitean overtones here.

\item $[7]$, \textit{we too will pursue} --- \ul{Direct statement about how this text will proceed}. Sextus will try to find contradictions within the disciplines in order to show that one particular stance can't be justifiably taken. Recall Plato's use of contradiction to show the unsatisfactory nature of etymological reasoning.

\item $[8]$, \textit{let us first look at the general} --- In general, Sextus will proceed from the general to the particular.

\end{itemize}

\section{General Arguments Against the Disciplines}

\begin{itemize}

\item $[9]$, \textit{dispute ... among philosophers} --- I.e., the nature of learning, whether learning \textit{per se} is possible (cf. Plato's \textit{Meno}).

\item $[9]$, \textit{there is not any discipline} --- Note that the problem is broken down into for analytical questions. The stakes of ``is not any'' is to show that the disciplines are incoherent, not that, in some trivial sense, people don't practice things labeled academic disciplines.

\item $[10]$, \textit{first about the first} --- The principle that `like attracts/resembles like' is fundamental to Greek thought.

\item $[10]$,\textit{taught by being or what-is-not-being} --- \ul{Fundamental methodological move for Sextus}, namely, dividing a phenomenon into opposites that, together, comprise the totality of the phenomenon. Then, if neither obtain, then the phenomenon doesn't obtain.

\item $[11]$ \textit{not possible for the same thing ... no attributes} --- Principles from Aristotle and Parmenides.

\item $[12]$, \textit{an appearance} --- Not clear exactly what this refers to, but the difference between appearance and reality is important, especially in Plato.

\item $[13]$, \textit{what is true is a being} --- \textit{Important assumption}.

\item $[16]$, \textit{nothing will be unteachable, ... nothing is teachable} --- Perhaps not immediately apparent, but keep in mind the preceding arguments about unteachable things. In general these arguments rest on the idea that something can't come from nothing, a problem that looms large in Greek metaphysics.

\item $[20]$, \textit{thing that are taught have to be sayables} --- Recall the separation of learning and names in \textit{Cratylus}.

\item $[23]$, \textit{see white ... smell fragrance} --- Note that Sextus was a physician of the `Empiricist` school.

\item $[25]$, \textit{the body composed out of them ... things themselves} --- Compositional nature of truth / being.

\item $[30]$, \textit{it is unclear} --- Perhaps a Platonic idea, that knowledge must be perfectly lucid. A perhaps unintuitive twist on it here, that contradictions within a discipline negate the entire discipline.

\item $[10-30]$ --- Schemas.

\qtreeinithook \qsetw \scriptsize
    \Tree [.[teachables] [.(epistemology) false [.true ] ] [.(ontology) what-is-not [.what-is not-something [.something [.incorporeal transition-from-experience [.experience ] ] [.body [.perceptible apparent ] [.intelligible unclear ] ] ] ] ] ]
    
\qtreeinithook \qsetw \scriptsize
    \Tree [.[teachables] non-expert [.expert experience [.speech non-signifying [.signifying fiat [.nature ] ] ] ] ]
    
\normalsize

\item $[37]$, \textit{by nature or by fiat} --- Cf. Plato's \textit{Cratylus}.

\item $[39]$, \textit{let's suppose}, Another common method for Sextus, assuming something as true for the sake of furthering his point. 

\end{itemize}

\section{Against the Grammarians}

\subsection{Introduction}

\begin{itemize}

\item $[41]$, \textit{from infancy} --- Literacy was extremely important to Greek culture and self-identity and was relatively high for its time period (though still low by modern standards).

\item $[42]$, \textit{human being by nature loves learning} --- Recall Plato's bid for rationality as humans' defining trait.

\item $[43]$, \textit{pronouncing rationally on the things in myths} --- Includes activity at the library of Alexandria.

\item $[44]$, \textit{letters ... complete grammar} --- Note that in Greek the words translated as ``literacy'' and as ``grammar'' are different modifications of the same word. Similar, somewhat, the English distinction between ``a letter (of the alphabet)'', `'letters (to grandma)'', and ``a person of letters (i.e., very literate)''.

\item $[44]$, \textit{Crates of Mallos, Aristophanes, and Aristarchus} --- Important scholars at the library of Alexandria (from before Sextus' time).

\item $[49]$, \textit{general expertise ... deeper power} --- Opposition between breadth and depth.

\item $[49]$, \textit{everyone agrees about that} --- \ul{No one disputes that reading is good, generally speaking. Question is whether grammatical expertise is useful.}

\item $[50]$, \textit{the goal of every expertise} --- \ul{Important assumption that expertise has a pragmatic goal, rather than some self-fulfilling purpose.} 

\item $[52]$, \textit{forgetfulness ... anything necessary is not possible} --- The pragmatic goal for literacy in the general sense. Plato argued against the benefit of writing w/r/t to memory in his \textit{Phaedrus}. This claim seems to assume learning in the sense of academic learning as in Alexandria (at least for, e.g., Medicine), as oral learning is likely sufficient for many necessary things, depending on how ``necessary'' is defined.

\item $[53]$, \textit{for if the lines of attack ... then literacy is useful} --- A condensed proof by contradiction, perhaps a tautological one insofar as it seems to assume literacy is good in the first place.

\item $[54]$, \textit{boastful and busybody} --- Slyly introducing pejorative terms for the second kind of grammar.

\item $[55]$, \textit{For the use ... observation of these} --- I.e., literacy has some impact on ethics and is in some sense traditional, i.e., based on interpersonal relationships.

\end{itemize}

\subsection{Scrutiny of prevailing conceptions of grammar}

\begin{itemize}

\item $[57]$, \textit{any consistent and real discipline} --- A common line of attack for Sextus. A lot of the ``definitions'' listed in this section are still relevant in modern conceptions of grammar/literacy/poetry.

\item $[57]$, \textit{Dionysius of Thrace} --- Scholar at Alexandria who wrote perhaps the first explicit grammar of the Greek language.

\item $[60]$, \textit{Ptolemy the Peripatetic} --- A peripatetic is someone who follows in the tradition of Aristotle. This person is only known through Sextus.

\item $[61]$, \textit{experience itself ... shared exercise} --- This definition of experience is somewhat similar to how Sextus will want to characterize the ``good'' grammar, though he refutes it here. Note again that it's interpersonal, and, surprisingly ``non-rational''(i.e., w/o reference to second-order thought?).

\item $[64]$, \textit{he would not say} --- Another proof by contradiction based on a binary, which assumes how the supposed interlocutor (Dionysius of Thrace, were he still alive) would respond. This is an important mode of argument even today, namely extending the implications of somewhat thought to reach conclusions they didn't explicitly make.

\item $[66]$, \textit{When they say ... infinite number} --- Important type of proof by contradiction, which ends in the impossibility of dealing with infinity in some way.

\item $[67]$, \textit{not of a few ... most of the things} --- I.e., a purely quantitative difference.

\item $[68]$, \textit{sorites} --- See the text's footnote on \textit{sorites} problems. The idea is that ``most'' is an undefinable amount.

\item $[68]$, \textit{question that goes little by little} --- The Socratic method, perhaps?

\item $[70]$, \textit{Just as ``fewer''... relation to ``few''} --- I.e., they are tautologically defined.

\item $[71]$, \textit{false that the grammarian knows} --- Note the turn here from meaninglessness to falsity. Because ``most'' is undefinable, ``most'' of something would be meaningless. But here, Sextus is saying you could legitimately call the grammarians' learning ``few'', when comparing it to infinity. (See the text's footnote.)

\item $[71]$, \textit{Asclepiades} --- of Myrlea, from Spain, wrote grammatical works.

\item $[73]$, \textit{grammarian, not of grammar} --- It's at least theoretically possible, if false, to define a thing apart from it's enactor (cf. \textit{Cratylus)}.

\item $[74]$, \textit{knowledge is nothing aside from the knower} --- Sextus refutes the above claim based on the idea that literacy is an activity of some sort.

\item $[75]$, \textit{apprehensions} --- See the text's footnote. Translates the important word \textit{techne}.

\item $[76]$, \textit{Chaeris} --- Grammarian, probably from Alexandria and pupil of Aristarchus.

\item $[77]$, \textit{method} -- Move from experience to method (still to be refuted).

\item $[79]$, \textit{the critic is different} --- Jakobson in the 20th century makes a distinction between the critic and the grammarian.

\item $[81]$, \textit{everything signified} --- I.e., a different infinity than the other definitions but still an infinity.

\item $[81]$, \textit{for the gods} --- Recall the special status/language of the gods in \textit{Cratylus}. I.e., there is, here, an infinity larger than the gods (or something along these lines).

\item $[82]$, \textit{science sets bonds} --- \ul{Important definition of Sextus' methodology.}

\item $[82]$, \textit{time likes changes ... one that is changing} --- The problem of time, also mentioned in \textit{Cratylus}. This infinity is also increasing. That there are different kinds of infinity is a fairly modern concept in the history of mathematics.

\item $[84]$, \textit{Demetrios} --- Little known grammarian.

\item $[85]$, \textit{things are said in the poets} --- Argument based on the content of texts.

\item $[88-89]$, \textit{all of them in general ... extend to every word} --- Not only does each language and sub-language proceed to an infinity, but grammatical abstractions can also proceed to an infinity.

\end{itemize}

\subsection{Transition from the conception of grammary to its content}

\begin{itemize}

\item $[91]$, \textit{one part of grammar ... special} --- \ul{Definition analysis of types of grammar}. Typical of Sextus' method to attack by breaking down into sub-categories.

\item $[92]$, \textit{set things in order} --- This is what Sextus will spend the most time on and is Sextus' own definition of grammar proper.

\item $[94]$, \textit{highly interwoven} --- Even Sextus' own outline is subject to lack of clarity.

\item $[96]$, \textit{if any one of these ... presence of the rest} --- Again note the method. Sextus will end up doing much ``more'', according to him, than is necessary, to show just how futile the discipline really is.

\end{itemize}

\subsection{The expert part of grammar}

\begin{itemize}

\item $[99]$, \textit{First in order ... unlettered} --- Note the method, which is to start with the smallest units (cf. \textit{Cratylus} and to show that the label of ``lettered'' is not only useless but false.

\item $[100]$, \textit{elements of the voice} --- Speaking takes precedence over writing.

\item $[104]$, \textit{composite of two} --- Cf. \textit{Cratylus} on secondary versus primary names.

\item $[108]$, \textit{what is receptive ... by a modulation} --- A perhaps not very convincing argument. At least, presumably some name for the letters could be concocted to make the analogy work.

\item $[110]$, \textit{if there is any ... but only short} --- A clearer/shorter summary of the preceding, more convincing than the previous, at least given certain metaphysical assumptions. Still, it could easily be posited that each ``two-timed'' letter is just two distinct letters, one long and one short, but Sextus will now attempt also refute this.

\item $[113]$, \textit{sons of grammarians ... other cases} --- These are basically the terms still used for Greek grammar, despite Sextus' arguments.

\item $[114]$, \textit{so that altogether} --- What is the argument here? That there are many elements? This doesn't seem to refute anything in a concrete way.

\item $[117]$, \textit{unchanging from start to finish} --- This is not the modern definition of a phoneme.

\item $[118]$, \textit{these too will be elements} --- Again the argument here is not clear. Why is it ``bad'' in general sense if certain grammarians are wrong? Why isn't Sextus simply doing ``better' grammar?

\item $[119]$, \textit{every element is said to be common} --- A more robust argument.

\item $[120]$, \textit{impasse} --- See the text's footnote on this point.

\item $[124]$, \textit{minimally short time} --- Linguistically speaking this is false, nor, even, is it true that all sounds considered ``short'', in a given language, need to be the same length. See the text's footnote on this, as well.

\item $[132]$, \textit{new methods} --- Compare to paragraph [53], which references teaching arguments. The tone is as if the text were a training manual of some sort for arguments / arguing. 

\item $[135]$, \textit{the discourse of which ... nothing} --- The argument here is not very clear. It seems to be that, if the discourse is simply tantamount to its parts, it doesn't make sense to speak of the discourse as a separate entity in the first place.

\item $[136]$, \textit{for if beyond ... of one another} --- Again, the motivations for this claim are obscure. What it would mean for a word to be ``part'' of another word is unclear, which is Sextus' point, but why he hypothesizes this in the first place as an option is unclear. Recall the compositional nature of sentences in \textit{Cratylus}.

\item $[141]$, \textit{we are in a position ... others too} --- Generally Sextus has continued on with explanations even when he says they are necessary. This is an instance of him curtailing his own method in order to focus on a representative type. Note that ``names'' are still the paradigmatic example.

\item $[144]$, \textit{How did grammatical ... either side} --- Probably a reference of some kind to \textit{Cratylus}. Note that this makes Plato a ``grammarian'' in some sense, rather than a philosopher. Sextus here, too, uses philosophical arguments as controlling or determining of grammatical arguments.

\item $[146]$, \textit{but if they say ... smoother} --- Sextus uses similar arguments to those in the \textit{Cratylus}, but he makes them about grammatical terms/functions rather than about word meaning. See the text's footnote for the collar metaphor.

\item $[149]$, \textit{This is not a reason ... and others the other} --- How a particular ``usage'' is ``determined'' is not clear here, but seems to rely on some idea of a community of speakers (foreshadowing, perhaps, Saussure). In any case, Sextus is making a bid for a conventional, non-natural basis for words.

\item $[153]$, \textit{But if the basis ... common usage} --- \ul{Sextus links naming/language to non-expert, i.e. common, usage}. In this way, claims to the naturalness of names is tied to (and refuted as) the domain of experts.

\item $[156]$, \textit{For how ... among philosophers} --- The argument here takes a different tack. Namely, instead of showing the analytical/inner contradiction of a certain definition, Sextus claims that because there is no agreement about what constitutes the definition of ``sayable'', it can't be valid for grammar. This is more directly a refutation of Stoic philosophy about what can be an ``incorporeal''. Sextus also seems to be implying that ``body'' is an ``incorporeal'', but this would be strange and not in line with his own earlier claims.

\item $[157]$, \textit{The person who says ... it was agreed} --- \ul{Important methodological point that people are required to give justifications based on agreed premises.}

\item $[163]$, \textit{everything that receives ... same state} --- It's not clear what the motivations for this claim are.

\item $[167]$, \textit{will only increase ... not the whole} --- The argument here for addition  seems to assume something here like Sextus claimed for subtraction in [163], namely, that every element in a set is affect by addition to that set.

\item $[170]$, \textit{futile} --- A reprise of the criterion of usefulness.

\item $[171-172]$, \textit{we should ... not useful} --- It's not clear what this ``bind'' would look like. Presumably the idea is that there should be some kind of ad hoc accord about which for to use until the debate is decided definitively for one. Instead, people seem to use whichever of a disputed form without incurring any problems.

\item $[176]$, \textit{there are two ... everyday talk} --- \ul{Important distinction between good and bad forms of criteria, analogy vs everyday usage.}

\item $[177]$, \textit{oblique cases of the nominative Zeus} --- Recall this being the first etymological example in\textit{Cratylus}

\item $[177]$, \textit{given before} --- Sections C/H.

\item $[178]$, \textit{Just as in a city ... nearly insane} --- The city analogy might reference Plato's \textit{Republic} and the ``insane'' person might refer to Hermogenes in \textit{Cratylus}. The word for ``private'' is the Greek \textit{idios} (from where we get ``idiot''), a traditionally very negative term in the city-state setting.

\item $[180]$, \textit{every expertise ought to take shape from some principle} --- \ul{Important methodological claim.}

\item $[180-]$,\textit{expert or non-expert} --- Back to Sextus' binary tree method.


\qtreeinithook \qsetw \scriptsize
    \Tree [.good-Greek non-expert [.expert self-justifying [.other-justifying usage [.expert^n clear [.unclear nature [.one-person/fiat assertion [.demonstration [.analogy some [.all ] ] [.practice in-tune [.out-of-tune ] ] ] ] ] ] ] ] ]
    
\normalsize

\item $[183]$, \textit{this criterion then ... need of expertise} --- This seems to imply that ordinary usage can, unlike expertise, take itself as its own principle/basis.

\item $[188]$, \textit{the many rather than the one} --- Recall that Plato explicitly wanted to deny the validity of word meaning being open to majority / democratic principles.

\item $[192]$, \textit{either all or most ... aware of it} --- Saussure will make much of quotidian analogy. Chomsky will make much of the fact that most people produce new utterances without knowing how.

\item $[194]$, \textit{clarity and smoothness} --- Two traditional ideals for Greek writing.

\item $[200]$, \textit{you are also rejecting analogy} --- Even reconstruction of obscure words requires analogy to some sense of usage. This neglects the potential usefulness of analogy for the reconstruction of lost forms.

\item $[206]$, \textit{laughed at} --- The negative version of usefulness for Sextus. Cf. [176]

\item $[208]$, \textit{just as Homer himself} --- In an ironic way, Sextus has, by rejecting Homeric usage, made his philosophy the most in line with Homer.

\item $[212]$, \textit{how was I in error ... commit a barbarism} --- Recall the very similar example toward the end of \textit{Cratylis}.

\item $[216]$, \textit{consistency} --- On the other hand, Plato also rejects ``consistency'' as a valid criterion in \textit{Cratylus}.

\item $[217]$, \textit{we do not say these things} --- That different verbs have different properties when given identical transformations is an important productive fact utilized in modern linguistics.

\item $[220]$, \textit{a barbarous one} --- Sextus poignantly turns the grammarians' terms against them.

\item $[221]$, \textit{certain universal theories} --- Critique here applies to Sextus' rival medical schools.

\item $[226]$, \textit{just as in many other cases in nature} --- This implicitly treats language as a natural phenomena, a relatively modern view.

\item $[227]$, \textit{we too} --- Cf. [191].

\item $[232]$, \textit{impasse-free} --- See the text's footnote here. Theoretically, the implication here, as later spelled out, is that the skeptic/Sextus is willing to adopt whatever usage is appropriate for any given context. But, strictly speaking, this implies that he would be fine adopting the ``usage'' of the grammarians, too.

\item $[234]$, \textit{same thing} --- What constitutes the ``same thing'', linguistically speaking has always been a central problem, including in \textit{Cratylus}.

\item $[235]$, \textit{with a view to those present} --- The community/usage base is more explicit here than before.

\item $[242]$, \textit{authentic names} --- Again, a central problem in \textit{Cratylus}, and Sextus recycles some arguments from it here. See the text's footnote.

\end{itemize}

\subsection{The special part of grammar}

\begin{itemize}

\item $[296]$, \textit{the the writers ... what is false} --- Sextus, like Plato, makes a bid for his own genre/mode as the right one. Poetry, though, unlike grammar, is also useful (cf. [370]). Note that the ``badness'' of grammar is heightened from previous sections. Now, grammar is not only useless, but false and misleading.

\item $[300]$, \textit{every piece ... objects or both} --- The referential nature of language is a large and never-ending debate in linguistics/philosophy. The exact nature of Sextus' sense of it is not explained here.

\item $[300]$, \textit{the grammarian is not at once all-wise} --- Remember that Sextus will go on to refuse several other disciplines.

\item $[301-302]$, \textit{where is ... Eudoxus} --- The argument here is very similar to one in Plato's \textit{Republic}.

\item $[303]$, \textit{like is known by like} --- Sextus' sense of the traditional import of this model is correct.

\item $[303]$, \textit{Empedocles call himself a god} --- Somehow this is supposed to make Empedocles seem not arrogant.

\item $[305]$, \textit{the sun ... blinds the sight} --- This is probably a direct reference to Plato's \textit{Republic}, Book 7, the famous ``cave allegory''.

\item $[311-312]$, \textit{what is in motion ... in the future} --- Note the return to the original style of binary argumentation, but in the voice of another philosopher.

\item $[314]$, \textit{even this is impossible ... we do not know} --- That there are an infinite amount of potential words wouldn't seem to mean all unknown words are unlearnable.

\item $[316]$, \textit{If only ... together} --- See the text's footnote on the poem.
 
\end{itemize}

\newpage
\part{Notes on Saussure's \textit{Course in General Linguistics}}

\section{Introduction}

\subsection{Chapter 1}

\begin{itemize}

\item \textbf{[13]}, \textit{it offers no scientific or objective ... forms} --- Note that the mode/goal of linguistic inquiry here is already different from Plato. Saussure wants to portray science as an end in itself, unlike Plato. Plato's and Sextus' sense of linguistics is in many ways, as Saussure denounces, interested most in ``forms''.

\item \textbf{[13]}, \textit{linguistic structure} --- This will be a key move for Saussure.

\item \textbf{[13]}, \textit{criticism} --- Note that Sextus Empiricus also makes a distinction between the grammarian and the critic, as does Jakobson after Saussure.

\item \textbf{[14]}, \textit{living language} --- Another key move.

\item \textbf{[15]}, \textit{these series of forms ... Sanskrit} --- Note that the comparative method requires criteria by which objective(-like) deductions can be made.

\item \textbf{[15]}, \textit{reconstruction} --- Word/form analysis has new, positive purpose compared to Plato and earlier grammarians: the reconstruction of unknown/unattested forms.

\item \textbf{[16]}, \textit{to define exactly what it was they were studying} --- This is crucial for the ``scientific'' basis Saussure wants to introduce, which requires a defined object with falsifiable properties.

\item \textbf{[17]}, \textit{a province of their own} --- Linguistics as another science, subject to the same criteria and part of the natural world and its laws. This is especially important for Chomsky later.

\item \textbf{[19]}, \textit{collective mind} --- Hinting at the social basis of language for Saussure, distinct from later cognitive theories of language like Chomsky's.

\end{itemize}

\subsection{Chapter 2}

\begin{itemize}

\item \textbf{[20]}, \textit{periods of decadence ... all forms of expression} --- Saussure's fighting against the idea that there are ``illegitimate'' versus ``pure'' forms/sources of certain languages.

\item \textbf{[20]}, \textit{the aims of linguistics} --- Focus on the ``laws'' of language, distinct from the ``generative''/cognitive approach of Chomsky.

\item \textbf{[21]}, \textit{distinguished from anthropology} --- Claude L\'{e}vi-Strauss famously uses results in Saussurian linguistics as input for anthropological claims.

\item \textbf{[21]}, \textit{the relation is unilateral} --- This isn't true in a certain sense. What phones (actual sounds) a person hears can at least be influenced by what phonemes (set of possible sounds in a given language) someone has available to them.

\end{itemize}

\subsection{Chapter 3}

\begin{itemize}

\item \textbf{[23]}, \textit{other sciences are provided} --- How distinct this picture of linguistic objects is from, e.g., the complex interactions of physics/chemistry/biology in cell functioning is perhaps overstated.

\item \textbf{[23-24]}, \textit{the ear perceives ... there is no way out of the circle} --- This list of difficulties sets up a fundamental (and very influential) methodological move for Saussure. Namely, linguistic phenomena is not unary in any of its domains but is, instead, a set of irreducible binaries. You might compare this, and the particular problem of language's origin, to a problem that Plato ran up against in \textit{Cratylus}, that of deriving an etymology from an irreducible, unanalyzable root/source.

\item \textbf{[25]}, \textit{one solution only ... only one part of language} --- This solution, that ``structure'' (a unary concept?), as a scientifically valid subset of linguistic phenomena, can account for binary relations becomes highly influential in a number of fields. For Saussure, linguistic structure manifests as a complete and socially defined set of linguistic entities. Chomsky will similarly reduce language to definable functional elements but will shift focus to what Saussure calls here the ``language faculty''. (Saussure's ``structure'' doesn't directly map onto this ``faculty''.) For Chomsky, this faculty exists in the brain and can generate the set of all possible strings in a given language.

\item \textbf{[25]}, \textit{faculty endowed by nature ... and conventional} --- The same debate invoked by Plato, whether linguistic phenomena is natural or conventional, but in a very different sense. Plato's question is ontological, whereas Saussure's is biological. While Saussure comes down on the side of convention, Plato's question has essentially no meaning for Saussure, since the latter has no conception of language as ``existing'' outside of its speakers. (Recall also that Plato explicitly disavows any socially based characterization of language.) Similarly, Chomsky comes down against conventionalism, but Saussure's question also is meaningless for Chomsky, since all (scientific) phenomena for Chomsky are ``natural''.

\item \textbf{[28]}, \textit{circuit} --- Recall that this book is culled from Saussure's lecture notes, so this drawing, and at least most of those to follow, would have been one reproduced by a student from something Saussure drew on a chalkboard.

\item \textbf{[28]}, \textit{facts of consciousness ... physiological process} --- This [set of abstractions] $\rightarrow$ [set of sounds] operation is, in schema, the model contemporary linguistics uses.

\item \textbf{[29]}, \textit{the sound patterns ... actual sounds} --- Similarly, this distinction is still very important and maps relatively well onto the contemporary distinction between phonemes and phones, respectively.

\item \textbf{[29]}, \textit{one must leave ... social crystallization} --- Saussure's social definition of linguistic structure, though the faculty (technically different here than that before, but probably ultimately referring to the same process) ``plays a major role''. The ``crystal'' metaphor common throughout 20th century theory, perhaps starting, or gaining relevance, with Marx.

\item \textbf{[30]}, \textit{distinguishing ... accidental} --- Another very important and influential distinction, between ``speech'' (French, \textit{la parole}), which is individual, and ``language'' (proper, itself French, \textit{la langue}), which is social. There is a third term that shows up here and there, usually also translated as ``language'' (French, \textit{le langage}), which is the superset of speech and language-itself.

\item \textbf{[31]}, \textit{certain ambiguous terms ... in other language} --- Note this is a different issue w/r/t to foreign language than that in Plato and Sextus. The latter were interested in the (supposed) fact that two words in different languages can have the same meaning/content.

\item \textbf{[31]}, \textit{contract} --- It might be historically interesting that Saussure uses this metaphor, as it is a central concept in political theory.

\item \textbf{[31]}, \textit{individual needs ... only gradually} --- The problem of learning a language is still important. Chomsky will stress, in fact, how quickly children seem to learn language, despite their getting little explicit input (known as the argument from the poverty of the stimulus).

\item \textbf{[31]}, \textit{dead languages ... linguistic structure} --- This is an important pragmatic distinction between Saussure and Chomsky. The latter will mostly disavow a method that depends on non-spoken language, since he refocuses linguistic inquiry onto the individual speaker.

\item \textbf{[32]}, \textit{dictionaries and grammars} --- Chomsky will strongly reject this.

\item \textbf{[33]}, \textit{semiology} --- Barthes and other literary theorists take up this task explicitly.

\end{itemize}

\subsection{Chapter 4}

\begin{itemize}

\item \textbf{[37]}, \textit{material substance} --- Note that, although signs in the language system are not ``abstract'', they are also not ``material''.

\item \textbf{[37]}, \textit{a language accumulates} --- An interesting statement as to how  language is / might be learned.

\item \textbf{[38]}, \textit{out of reach ... individuals} --- This is how Saussure gets out of Socrates' attack on the utter uselessness of language that is purely conventional.

\item \textbf{[38]}, \textit{it would be impossible ... same point of view} --- Recall that, perhaps contra this statement, the structure of language was invoked (p. [25]) to solve the problem of linguistic phenomena being constituted by irreducible binaries. 

\end{itemize}

\subsection{Chapter 5}

\begin{itemize}

\item \textbf{[42]}, \textit{remote valleys ... as examples} --- Key claim if one is to distinguish language as a human phenomenon as opposed to a set of historical phenomena.

\item \textbf{[42]}, \textit{a borrowed word ... context of a system} --- Loan words do in fact remain marked as ``foreign'' for long periods after their initial borrowing.

\end{itemize}

\subsection{Chapter 6}

\begin{itemize}

\item \textbf{[44]}, \textit{actual object ... comparing them} --- A fairly contemporary sounding statement as to the task of descriptive linguistics.

\item \textbf{[44]}, \textit{impossible to ignore ...represented} --- Again, Chomsky will generally assume that writing can be ignored.

\item \textbf{[48]}, \textit{the tendency ... for an idea} --- Note that Saussure didn't speak any Chinese dialect, so this claim is pure, likely unwarranted, speculation.

\item \textbf{[52]}, \textit{writing is not a garment, but a disguise} --- Recall the longstanding relationship between ``texts'' and ``textiles'', as seen especially in Callimachus and in Plato's weaving examples.

\item \textbf{[53]}, \textit{pronunciation ... by its history ... the only defensible pronunciation} --- Saussure's system might in fact reject this, claiming instead that (phonemic) pronunciation is determined by the linguistic structure. The structure provides the phonemic rules, so ``what people say'' should be the criteria for which is correct. This comment uses Sextus Empiricus' critique of formalism, as well as Chomsky's.

\item \textbf{[54]}, \textit{sept femmes ... vingt} --- In fact the `t' in `sept' and `vingt' is pronounced now in French, but the `g' in `vingt' is still not pronounced.

\end{itemize}

\subsection{Appendix}

[skip]

\section{Part 1}

\subsection{Chapter 1}

\begin{itemize}

\item \textbf{[98]}, \textit{linguistic sign ... actually a sound} --- What actually constitutes a ``sound pattern'' (French ``image acoustique'', literally, ``acoustic image'') is not always clear in Saussure, and in his writings the word that refers to this part of Saussure's theory changes several times. Consistent is the idea that the phonic part of the sign is not ``physical''. The text's footnote here is helpful.

\item \textbf{[99]}, \textit{linguistic sign ... following diagram} --- Note the text has a typo, ``p. 67'' instead of ``p. 77''. The arrows in this diagram seem to indicate that both parts co-refer. In other writings Saussure refers to the dual nature of the sign as being two sides of a coin (or paper, later in this text (p. [157])), which may be what this diagram is trying to portray, while still others portray a rectangle bisected diagonally (i.e., into triangles) to display the two parts.

\item \textbf{[99-100]}, \textit{we propose ... be replaced} --- This is a fundamental three-part distinction for Saussure: 1) signification, French \textit{le signifi\'{e}}, the thing signified, the concept; 2) signal (signifier), French \textit{le signifiant}, the thing signifying, the sound pattern; 3) sign, French \textit{le signe}, signified+signifier, the totality. This division, a subset of the linguistic structure (\textit{le langue}, language-itself) conceptually maps well onto the larger three-partite division between: speech / language-itself / language (superset). Recall that Plato also had a three-partite division: the thing named, the name, and the (act of) naming. Whether Plato's act of naming maps onto Saussure's sign is unlikely but debatable.

\item \textbf{[100]}, \textit{arbitrary} --- Saussure's notion of the arbitrary in no way means `random' or `capricious', like some connotations in English. It simply means that there is no rational or natural explanation for a given signifier-signified link.

\item \textbf{[100]}, \textit{no one disputes ... arbitrary} --- Many linguists and theorists after Saussure attribute the concept of the sign's arbitrary nature to Saussure himself. This is demonstrably historically false, as even Saussure suggests. The concept of arbitrariness as a (potential) linguistic theory goes back, as we've seen, at least to Plato (i.e. ``convention''), and such a theory was espoused by Aristotle.

\item \textbf{[102]}, \textit{compare a French ... originally meaningful words} --- Note English `bow-wow' is related here. Also, recall the exclamation \textit{hie, hie} from Callimachus, which might be related.

\item \textbf{[103]}, \textit{certain temporal characteristics ... incalculable} --- Though Saussure warns against it, the linear nature of linguistic phenomena is often neglected by followers of Saussure, though not by contemporary linguistics.

\end{itemize}

\subsection{Chapter 2}

\begin{itemize}

\item \textbf{[104]}, \textit{from the point of view ... chosen} --- Again, this is how Saussure ensures that a set of conventional signs does not amount to absolute uselessness. (I.e., the danger, at least for Plato, is that a set of purely conventional signs would create a situation in which nobody knows everybody else's word for anything.)

\item \textbf{[104]}, \textit{contract} --- Recall Saussure's prior, opposite-valence use of this word (p. [31]).

\item \textbf{[105]}, \textit{the initial assignment of names ... taking place} --- Perhaps a direct reference to Plato's \textit{Cratylus}.  Saussure relegates to triviality what was a fundamental problem for Plato, the initial origin/makeup of words. This is in part due to the ontological versus biological/social framework that splits the two theorists.

\item \textbf{[105]}, \textit{no society ... accept} --- Though Saussure doesn't present it as such, this claim is tautological given Saussure's social definition of linguistic structure.

\item \textbf{[106]}, \textit{the continuous efforts} --- Chomsky will suggest that a language is very quickly learned by children, and modern psycholinguistics confirms that the ability to learn language is lost after a certain age. Still, Chomsky would agree that there are no well-demarcated linguistic ``generations''.

\item \textbf{[106]}, \textit{largely unaware} --- Chomsky will also make much of this fact.

\item \textbf{[106]}, \textit{the arbitrary nature ... to change it} --- The formal distinction here is unclear, though it seems to fall along prescriptive versus descriptive lines. The argument could be reconstructed as referring to different valences of the arbitrariness: the sign could change because there's no ``right'' sign but it doesn't change because, by extension, there is no ``wrong'' sign.

\item \textbf{[107]}, \textit{the great number ... influence upon it} --- These reasons might be thought of as referring to the finite nature of human cognition. There may be theoretical worth in linking these reasons, and that language exists in time, more explicitly to the linear character of the signifier itself.

\item \textbf{[108-109]}, \textit{the final analysis ... principle of continuity} --- Probably influenced here by the Neogrammarians.

\item \textbf{[109]}, \textit{the factors ... signal and signification} --- A response to and advancement from the Neogrammarians.

\item \textbf{[110]}, \textit{any sequence of sounds} --- Of course, the human vocal tract cannot produce the entire, infinite set of conceivable of sounds.

\item \textbf{[111]}, \textit{Esperanto} --- Among the more influential of constructed languages. One of Saussure's younger brothers, Ren\'{e}, was an important leader in the Esperanto ``movement''. 

\item \textbf{[111]}, \textit{hen hatching a duck's egg} --- This is a bad analogy because Saussure here is not making an a priori morphological (biologically speaking) claim about language, as the analogy would suggest, but is instead making an a posteriori claim about language (as he soon will admit). 

\item \textbf{[112]}, \textit{avoiding ... himself understood} --- To put a finer point on Saussure's argument, one might point to the fact that a phoneme can be realized as a diverse set of phones.

\item \textbf{[112-113]}, [diagrams] --- These diagrams seem to me relatively incoherent, and the text is clearer in this case.

\item \textbf{[113]}, \textit{if a language ... no change to observe} --- The course notes agree on this claim, but it makes little sense, even in Saussure's system. A fitting entailment here is that an individual living in isolation would produce a large, perhaps infinite, amount of linguistic changes, since there would be no social constraints. Saussure himself admits that phonetic modifications first occur in speech, not in the structure. I.e., this is the case in which time is factored but society is not. The next case, on the other hand, is that in which society is factored but time is not, and this situation (not the former, as the text reads) would produce no linguistic change. In any case, Saussure's point is that time and society are contrastive forces that produce a constrained amount of linguistic change.

\end{itemize}

\subsection{Chapter 3}

\begin{itemize}

\item \textbf{[114]}, \textit{very few linguistics ... approaches} --- Despite Saussure's warning, much contemporary use of Saussure absolutely ignores the ``Diachronic'' portion of this work.

\item \textbf{[114]}, \textit{economics ... duality} --- Economics in its contemporary form, like linguistics at the time, was undergoing a series of substantial of changes not unrelated to the series of massive political changes then going on in Europe.

\item \textbf{[115]}, \textit{value} --- The concept of ``value'' is one which saw major theoretical upheaval.

\item \textbf{[116]}, \textit{pure values} --- What constitutes linguistic ``value'' for Saussure is not clear at this point but will be explained in Part 2. Whether there is such a sharp distinction between value in economics is debatable.

\item \textbf{[117]}, \textit{likewise ... evolution} --- A crucial terminological distinction.

\item \textbf{[117]}, \textit{Only by ... language user} --- This is a crucial distinction between Saussure and Chomsky, of which their difference on the use of written materials is a subset. Chomsky will stress that only native speaker intuitions are methodologically valid for fully accessing a particular language's structure. Saussure suggests here that a linguist can in principle fully access the structure of any given language.

\item \textbf{[118]}, \textit{traditional grammar ... observing facts} --- The emphasis on descriptive rather than prescriptive methods is still paramount in linguistics.

\item \textbf{[119]}, \textit{at the right moment} --- The point here is that the linguistic phenomena in question are completely fortuitous and are in no way necessary.

\item \textbf{[120]}, \textit{Another example ... geese} --- These original Germanic declensions described here, before the sound changes, are, interestingly similar to a set of declensions in Latin and Classical Greek.

\item \textbf{[121]}, \textit{the diachronic developments ... plurals of nouns} --- This must be the case for the sign to remain arbitrary.

\item \textbf{[122]}, \textit{at each stage ... brings it to life} --- The scientific meaning of this metaphor, even in Saussure's thought, is unclear, though its biblical source is dripping. Presumably it means something like ``gets used by society''.

\item \textbf{[122]}, \textit{for the changes ... coexisting terms} --- Why ``unintentional'' is contrasted to ``significant'' (the translation is literal) is obscure to me. Presumably Saussure doesn't mean that synchronic terms are somehow ``intentional''.

\item \textbf{[122]}, \textit{in the diachronic ... affected by them} --- This seems overstated. Diachronic changes are likely limited by synchronic facts to a finite set of possibilities. If not, the differential nature of the synchronic system (which Saussure discusses later) would be incoherent.

\item \textbf{[124]}, \textit{we can recognize ... the sign zero} --- The introduction of null operators is extremely methodologically important for linguistics, historically and currently.

\item \textbf{[125]}, \textit{game of chess} --- ``Chess'' was and becomes increasingly through the 20th century an important conceptual metaphor for a lot theoretical and artistic work. The broader metaphor/concept of (discrete) ``games'' is longstanding, even in Plato, but these more frequently fall under the subcategory of ``games of chance''.

\item \textbf{[126]}, \textit{unchanging principles of semiology} --- Note that these haven't been defined by Saussure, nor will they. Their general nature is in fact obscure here but the search for such general constraints on language is a driving force in contemporary linguistics.

\item \textbf{[127]}, \textit{there is only one ... upon the system} --- One might ask whether or not the advent of computer chess algorithms fixes this problem.

\item \textbf{[128]}, \textit{the one and only reality} --- This is difficult to parse. Previously Saussure had based the societal constraint on change (p. [107]) on the inability for any speaker to cognize the linguistic system in its entirety. In what sense, then, except by tautology, the synchronic system is ``real'' and the diachronic not ``real'' is at best unclear. Saussure had, in fact, just previously (p. [123]) claimed that the Latin stress system is in some sense operant in modern French, if differently than the modern French system proper.

\item \textbf{[128]}, \textit{object of synchronic ... sub-dialects} --- The attention to ``microvariation'' is still important to current linguistics.

\item \textbf{[130]}, \textit{ghost} --- Marx, prior to Saussure, will also describe economic value as having phantom-like properties.

\item \textbf{[133]}, \textit{real question ... material nature} --- Recall that Plato in \textit{Cratylus} made a sharp distinction between the ``elements'' or words and the words themselves. This becomes a serious problem for Plato in \textit{Theaetetus}, which seems to come out here, namely: How can a whole (a word) be, in some sense, more than its parts (elements/sounds)? See the translator's footnote here, too.

\item \textbf{[134]}, \textit{to summarize ... general about them} --- Perhaps this maps in some sense onto the descriptive (synchronic) versus prescriptive (diachronic) divide, though Saussure would posit all linguistics as a descriptive science.

\item \textbf{[138]}, \textit{the primitive stem} --- It is actually debated whether this is the correct original stem. The phenomenon in question here is known as Grassmann's Law.

\item \textbf{[138]}, \textit{studying languages and studying speech} --- The translation here should be singular ``language'' not plural ``languages'', as it translates the singular \textit{langue}, the ``linguistic system''.

\item \textbf{[138]}, \textit{analogical} --- The concept of ``analogy'' will become important later.

\item \textbf{[139]}, \textit{incorporation} --- The precise mechanism/definition for how certain diachronic changes become ``accepted'' is unclear throughout Saussure.

\item \textbf{[139]}, [diagram] --- Again the plural ``languages'' in this diagram translates the singular \textit{langue}, the linguistic system. While the singular ``language'' translates the singular \textit{langage} (see p. [31]).

\item \textbf{[139]}, \textit{these different ... constant principles} --- This is a very important move, namely that all languages, despite appearing different, have at least some shared properties. Chomsky will bring this assumption to its limits.

\end{itemize}

\section{Part 2}

\subsection{Chapter 1}

\begin{itemize}

\item \textbf{[142]}, \textit{unimportant changes} --- How to demarcate important from unimportant changes is unclear.

\item \textbf{[143]}, \textit{no demonstration ... of the data} --- This methodological move is consistent through Chomsky, as well.

\end{itemize}

\subsection{Chapter 2}

\begin{itemize}

\item \textbf{[144]}, \textit{signs comprising ... of that science} --- Any (in broad terms) science has particular ``objects'' of inquiry, so this isn't unique to Saussure. Modern linguistics, however, while having ``objects'' of study, will not, like Saussure's approach, be interested in ``entities'' but instead will be interested in structures and/or algorithms.

\item \textbf{[144]}, \textit{we are then left ... concrete object} --- This is an important distinction. In ``practice''/reality, the sign is not analyzable but can be analyzed in terms of abstractions. This solves Plato's problem of deriving a distinction between an atomistic entity and what appear to be its elements, namely by categorizing the two differently (real versus abstract).

\item \textbf{[144]}, \textit{the same is true ... signal} --- See the translator's footnote to this sentence, which is ludicrous. Whether it's justifiable to split sound from meaning for linguistic/analytic purposes is debatable. (Saussure says no, Chomsky says yes.) It's easily conceivable how meanings might be studied apart from sounds, and there is, in fact, an entire subfield of linguistics (semantics) devoted to this very pursuit.

\item \textbf{[145]}, \textit{a linguistic entity ... mechanism of the language} --- The implication seems to be that we don't have explicit access to non-delimited, i.e., non-``unit'', entities.

\item \textbf{[145]}, \textit{when we listen ... into pieces} --- Perhaps this foreshadows Chomsky's criterion of native speaker intuition. Though, again, Saussure denies one must be a native speaker for adequate analysis.

\item \textbf{[146]}, \textit{the unit has no ... a certain concept} --- Note that, technically, this definition includes units of any size, from subwords (morphemes) to words to superwords (phrases/sentences).

\item \textbf{[146]}, \textit{anyone familiar ... parallel sequences} --- Since speech is not the linguistic system itself, theoretically this method will entail gaining evidence from more than one speaker.

\item \textbf{[147]}, \textit{compare series ... by the sense} --- This is a risky methodology, insofar as, as modern linguistics has shown, many superficially similar looking strings can be derived from very different underlying structures. Saussure goes on to admit this only up to apparent homonyms, but the error can be more subtle than this.

\item \textbf{[147]}, \textit{what a word is ... concrete unit} --- Recent linguistics will similarly deprivilege the ``word'' as a primitive unit.

\item \textbf{[148]}, \textit{or else dispense ... concrete unit} --- The difficulty here, though not in all cases, seems fairly easily solved with a concept Saussure has already introduced, namely: null operators that attach to stable stems. Then, both the singular and plural similarity could be maintained, and the subword units could be applied to both cases. Of course, the word itself would still not be the domain of ``units'' proper.

\item \textbf{[148]}, \textit{lack of resemblance} --- Chomsky also stresses this.

\end{itemize}

\subsection{Chapter 3}

\begin{itemize}

\item \textbf{[150]}, \textit{the converse does not hold} --- The takeaway here seems to be that proven (methodologically checked) identities are a subset of all possible identities.

\item \textbf{[151]}, \textit{two trains} --- Chomsky has a very similar example of the persistence of certain  concepts through otherwise major changes, using a river instead of a train.

\item \textbf{[152]}, \textit{reality} --- Recall that ``reality'' doesn't necessarily mean ``material reality''.

\item \textbf{[153]}, \textit{conclusion is that ... linguistic reality} --- Modern linguistics will essentially explode traditional parts of speech into abstract analyzable units, as Saussure suggests here might be necessary.

\item \textbf{[153]}, \textit{concrete linguistic entities ... linguistic reality} --- Of course, Sextus Empiricus would reject this out of hand.

\item \textbf{[154]}, \textit{to determine units ... poorly defined units} --- As it's defined here, this is literally an infinitely large task, since it includes units larger than words, of which any given language has an infinite amount.

\end{itemize}

\subsection{Chapter 4}

\begin{itemize}

\item \textbf{[155]}, [summary] --- The concept of ``value'' will be described further in the following pages. For now, the following summary can be made:
    \begin{itemize}
    \item entity $\rightarrow$ concrete, real (inaccessible) object of linguistic structure
    \item unity $\rightarrow$ a well-defined, accessible entity
    \item identity $\rightarrow$ a set of unities deemed to be ``the same''
    \item reality $\rightarrow$ (still undefined) classes of unities
    \item value $\rightarrow$ a criterion of identity, a unit's differential scope w/r/t other identities
    \end{itemize}

\item \textbf{[155]}, \textit{psychologically ... shapeless mass} --- Of course, the empirical basis for this is unclear.

\item \textbf{[155]}, \textit{philosophers  ... constant way} --- And, of course, this is historically false. As we've seen Plato in \textit{Cratylus} hints at the possibilities of thought without language.

\item \textbf{[156]}, \textit{series of adjoining ... segmentation} --- This aspect of Saussure's theory seems fairly unjustified. It's not clear why, if sound and thought are divisible together, why the wouldn't in principle be divisible separately. If it's that each side need an outside delimiter, then an obvious ordering problem arises, as neither one could occur first nor could they occur at the same time. This could be an empirical claim, but Saussure certainly doesn't present it as such, himself calling it a ``mysterious process''.

\item \textbf{[157]}, \textit{a form, not substance} --- Hence, they can be real but not material.

\item \textbf{[157]}, \textit{values have no other ... agreement} --- By Saussure's own admission this doesn't seem to be the case, since, theoretically, it would still be useful, even for an individual, to order one's thoughts.

\item \textbf{[159]}, \textit{values always involve ... consideration} --- These features seem drawn from economics, as the later examples confirm, continuing Saussure's other references to the field.

\item \textbf{[160]}, \textit{a word can ... another word} --- Theoretically the ``dissimilar'' relation (between sound and concept) must come first, because otherwise the ``similar'' relation (between words) would be all muddled and indistinct, a ``mass''. This, unlike above, doesn't create an ordering problem, since Saussure is at pains to distinguish ``meaning'' from ``value''.

\item \textbf{[160]}, \textit{set of synonyms} --- In English, too, there might be cases in which ``to fear'' `works better' than ``to be afraid'' on account of their relative meaning.

\item \textbf{[161]}, \textit{if words had ... language and another} --- This might be taken as an empirical claim as to the ordering problem discussed above (p. [156]). Modern linguistics will attempt, however, to find linguistic universals, though these aren't concepts per se.

\item \textbf{[162]}, \textit{purely differential} --- A very important and highly influential formulation.

\item \textbf{[163]}, \textit{no linguistic ... with the rest} --- This would likely predict that there are an infinite amount of possible sounds, which Saussure is fine with (cf. p. [110]).

\item \textbf{[163]}, \textit{two signs a and b ... between a and b} --- It's not clear whether this contradicts the necessity for the sign to be established before its value.

\item \textbf{[164]}, \textit{sound is merely ... language uses} --- Recall Plato's analogy for letters/elements to paints in painting. 

\item \textbf{[166]}, \textit{their combination is a fact of a positive nature} --- This is important if, again, the sign has to be delimited before differential values can hold.

\item \textbf{[167]}, \textit{a unit ... differential in nature} --- ``Unit'' here seems to mean a set of one or more signs.

\end{itemize}

\subsection{Chapter 5}

\begin{itemize}

\item \textbf{[171]}, \textit{connexion in the brain} --- Saussure seems to mean this as something separate from the differential nature of values in the linguistic system. It appears that this, as well as the syntagmatic function, refers to a ranked set of differential values. At least the associative function seems idiosyncratic to individual speakers. Both functions seem here to refer to speech, rather than linguistic structure proper, and to be relevant only during a particular speech event. Saussure will go on to deny this, perhaps not convincingly.

\item \textbf{[173]}, \textit{constructed on regular patterns} --- The problem here is to figure out the scope of ``regular patterns''. Any utterance could, given a general enough definition of pattern, be attributed to some pattern. Still, the claim here is interesting, in that it posits that language has new forms ``built in'' and ready to use before they are used. Chomsky will make a similar claim.

\item \textbf{[174]}, \textit{impossible to say ... what order} --- Again, how much this is up to an individual brain and how much to the language system itself is unclear.

\end{itemize}

\subsection{Chapter 6}

\begin{itemize}

\item \textbf{[176]}, \textit{groups of both kinds ... language functions} --- Again, the question is how far to take this. In the case of \textit{-eux}, below, perhaps it is clear. However, theoretically any noun could come after, e.g., the word \textit{is}, so Saussure would have to strain to make this a part of structure in his conception.

\item \textbf{[177]}, \textit{the whole depends on the parts} --- It's potentially nonsensical to speak about ``parts'' and ``wholes'', which are nonlinear concepts, in the context of a linear relation. 

\item \textbf{[179]}, \textit{our memory holds ... kind and length} --- This is, as it stands, literally impossible, given that our brains have finite memory capacity.

\item \textbf{[181]}, \textit{not all signs are absolutely arbitrary} --- Saussure doesn't seem to need this concession to the arbitrary nature of the sign. He has already concluded that signs do not map directly onto words, but the phenomenon from which he is defining relative arbitrariness is compound words. He could have, however, simply derived compound words via the syntagmatic/linear process that applies to any potential sequence of words. In that case, compound words have no special status, but are an automatic feature of syntagmatic operations.

\item \textbf{[183]}, \textit{lexicalogical ... grammatical} --- This distinction seems to map onto the modern (descriptive, pre-theoretical)  distinction between analytic and synthetic languages, respectively.

\end{itemize}

\subsection{Chapter 7}

\begin{itemize}

\item \textbf{[186]}, \textit{linguistically, morphology ... distinct from syntax} --- This assumption has relatively recently come to be re-accepted in modern linguistics.

\end{itemize}

\subsection{Chapter 8}

\begin{itemize}

\item \textbf{[189-190]}, \textit{to associate ... notion of parts of speech} --- The terminology here seems to be confused. Probably ``real'' or ``existing'' or, as later, ``concrete'' is meant for ``material''. ``Abstract'', previously denied as a term for linguistic entities proper, seems to refer to classes of linguistic entities.

\item \textbf{[191]}, \textit{word order is ... abstract entity} --- The exact meaning/import of this is unclear.

\end{itemize}

\section{Part 3}

\subsection{Chapter 1}

\begin{itemize}

\item \textbf{[194]}, \textit{historical phonetics ... descriptive phonetics} --- The French text reads only ``phonetics'', the ``historical'' is the translator's addition. Similarly, ``descriptive phonetics'' should be translated ``phonology''. The phonetics-phonology distinction is as important to Saussure (as he has stressed) as it is to modern linguistics.

\item \textbf{[194]}, \textit{if the evolution ... synchronic with grammatical} --- Recall that, for Saussure, sound change entails meaning change, insofar as the value of both sides of the operation are differential. Hence, the series of complications that are to follow. 

\item \textbf{[196]}, \textit{we can say ... German perfect} --- Saussure doesn't need to admit this. He has already claimed that the differential systems of separate languages have different effects. What it would mean to synchronically compare two differential systems is not clear.

\end{itemize}

\subsection{Chapter 2}

\begin{itemize}

\item \textbf{[199]}, \textit{restricted by specific conditions} --- That sound changes are dependent on their context is still an important empirical fact in contemporary phonology.

\item \textbf{[199]}, \textit{spontaneous and combinative} --- The distinction here seems largely heuristic, rather than theoretically necessary. Namely, it seems to be a way to judge between two or more possible options for explaining a particular phonological change/law. The methodological examples to follow confirm this.

\item \textbf{[202]}, \textit{although it is not ... change occurred} --- This is a fair example of how empirical linguistic research currently proceeds. Namely, there is a disparate, often contradictory-seeming, set of linguistic facts that need to be explained in the most succinct way possible.

\item \textbf{[202-203]}, \textit{racial predispositions ... in Ionic} --- This is a relatively progressive stance for someone at the time. Saussure's other younger brother, L\'{e}opold, who worked in the French navy and studied Chinese astronomy, was a vocal proponent of scientific racism.

\item \textbf{[204]}, \textit{the difficulty ... opposite happens} --- Modern linguistics will portray this difficulty as one between two contrastive forces in phonological phenomena. On the one hand, there is a goal for minimum physiological effort; on the other hand, there is a goal for maximum phonological precision.  

\item \textbf{[204]}, \textit{scarcely possible ...difficult to pronounce} --- This is likely empirically false. Certain sounds are empirically more difficult to say than others, which explains the relative paucity of certain sounds (e.g., the ``th'' sound in English) across all languages

\item \textbf{[205]}, \textit{the many hesitations ... around him} --- This might be taken as a version of Chomsky's argument from the poverty of the stimulus.

\item \textbf{[206]}, \textit{act intermittently} --- Recall that Saussure defined phonological change as an ``imperative law'' (p. [134]).

\item \textbf{[208]}, \textit{laws of imitation} --- The problem of ``imitation'' is perennial in (Western) philosophy since (at least) Plato.

\item \textbf{[209]}, \textit{if grammar entered ... phonetic evolution} --- Apart from the grammatical consequences in the next chapter, Saussure just previously admitted that grammatical/conceptual changes can occur during political upheavals. What must be the case, here, then, if Saussure is to be consistent, is that sound and grammatical changes can occur independently and can subsequently influence each other, but they can't each be criteria of change for the other.

\end{itemize}

\subsection{Chapter 3}

\begin{itemize}

\item \textbf{[213]}, \textit{Xenophon's contemporaries ... were vocalic} --- This is how it's taught in current classics textbooks, in fact.

\item \textbf{[217]}, \textit{overlooking a synchronic duality} --- I.e., it's still important to take account of synchronic facts when doing diachronic analysis.

\item \textbf{[219]}, \textit{these sound contrasts ... synchronic opposition} --- The idea here is that particularly noticeable/definable oppositions are not different in kind than any other synchronic opposition. What this means, as implied before, is that synchronic oppositions have no need to be in any sense parsimonious. 

\end{itemize}

\subsection{Chapter 4}

\begin{itemize}

\item \textbf{[221]}, \textit{analogy} --- This is, in my view, an extremely important/elucidating chapter, but one that has often been neglected in literary and cultural studies.

\item \textbf{[221]}, \textit{analogy presupposes ... fixed rule} --- Again, the concept/function of imitation is highly fraught, as Saussure just admitted (p. [208]). Furthermore, the exact theoretical notion of analogy-as-image is muddled and changes throughout Saussure's lectures, on which the current text is based, as well as his other writings. I.e., it will remain relatively unclear what the precise conceptual motivations for analogy are. In any case, what can be said is that analogy is some interplay between a language's structural possibilities and which of those possibilities are, by some set of individuals brains, at some point actually produced.

\item \textbf{[222]}, \textit{in Greek, the perfect} --- The classical Greek perfect system, which is highly idiosyncratic, has long been the source of much linguistic analysis.

\item \textbf{[224]}, \textit{every analogy is a drama ... itself} --- This strikes me as a fairly odd and, for Saussure, unique analogical (both using and about ``analogy'') example. Drama has been the center of discussions on imitation since Plato, and three was the standard number of actors in classical Greek tragedy.

\item \textbf{[224]}, \textit{since language dislikes ... disappears} --- Of course, ``dislikes'' is in no way a technical or even precise term. The mechanisms which turn an alternate, analogical form into the default form and which delete the prior term are ambiguous throughout Saussure. This mechanism is important because presumably it's a question of when a particular word becomes explicitly part of the linguistic structure as a whole.

\item \textbf{[226]}, \textit{creations} --- Saussure is setting himself up for the implication that analogical creations are potentially limitlessly producible. The implication is that all analogical productions are inherent, i.e. potentially available, to a language at any given time. This would be a synchronic claim (made explicitly shortly); however, analogy is curiously in the diachronic section and has some relationship to time, hence the importance of the mechanism described above, still unexplained.

\item \textbf{[227]}, \textit{two things must ... belongs to speech} --- This still paints the system as prior to the individual speech, which gets confirmed in the following discussion. This clarification still doesn't explain how the analogical form comes to be, in aggregate, accepted as the default.

\item \textbf{[227-228]}, \textit{to summarise ... synchronic} --- Why is this concept in the diachronic section? In fact, several of the course lecture notes imply that analogy is much more diachronic than this firm statement would lead one to believe.

\end{itemize}

\subsection{Chapter 5}

\begin{itemize}

\item \textbf{[232]}, \textit{language retains ... vocabulary and grammar} --- The criteria for analogical acceptance is ``numerous'' enough. This, of course, is in no way precise. If analogy operates like other diachronic laws, at least in some sense, then this is not a unique ambiguity. If analogy is strictly a synchronic affair, this ambiguity poses a problem for Saussure.

\item \textbf{[233]}, \textit{the decisive argument} --- Here analogy seems to be a subcase of the verification process from before (p. [147]).

\item \textbf{[234]}, \textit{the action the language takes ... fixed system of units} --- This seems like another place where the synchronic definition of analogy is having trouble. What sense does it make to claim that a self-complete language system, at any fixed point, ``hesitates''?

\item \textbf{[235]}, \textit{reflects the process of change} --- This is how Saussure tries to get out of the synchronic-diachronic problem, but ``reflects'', like ``imitation'' and ``dislikes'', is not theoretically precise.

\item \textbf{[235]}, \textit{language is a dress} --- Another text-textile reference (cf. p. [52], ``language is not a garment, but a disguise''). 

\item \textbf{[237]}, \textit{the language is seeking} --- Oddly sentience-implying terms.

\end{itemize}

\subsection{Chapter 6}

\begin{itemize}

\item \textbf{[238]}, \textit{popular etymology} --- Many of the derivations in \textit{Cratylus} would fall under this category.

\item \textbf{[238]}, \textit{only difference ... howlers} --- This difference could be relabeled as `productive' and `unproductive', using contemporary terms. 

\item \textbf{[239]}, \textit{there is the case ... French suffix} --- Cf. the etymology of, e.g., ``island'' for an example of this in English.

\item \textbf{[240]}, \textit{analogy always presupposes ... rival to appear} --- This seems to contradict the process of analogy describe earlier (chapter 4), in which the old and new form can, for a time, coexist.

\end{itemize}

\subsection{Chapter 7}

\begin{itemize}

\item \textbf{[242]}, \textit{agglutination} --- This will prove to be a different case than that mentioned before w/r/t whether syntagmatic relations are sufficient for dealing with word combinations (cf. p. [181]). Syntagmatic relations proper could likely handle some of these cases on their own. But, for Saussure, what's crucial here, unlike the former cases, is that the agglutinated unities are perceived as such and can therefore themselves uniquely operate in the differential system.

\item \textbf{[244]}, \textit{intention} --- While analogy has previously been categorized by Saussure as ``psychological'', it has not been called ``intentional''. To attribute ``intention'' to these processes seems dangerous (cf. p. [122], though).

\end{itemize}

\subsection{Chapter 8}

\begin{itemize}

\item \textbf{[247]}, \textit{the following examples ...as a prefix} --- The example is the derivation of prepositions from adverbs. These two categories are still often semantically and grammatically linked, as well as difficult to disambiguate, e.g., the word \textit{before}. In \textit{I went before to the market} it's an adverb, in \textit{I went before him to the market} it's a preposition, and in \textit{I went before I went to the market} it's a conjunction. Modern linguistics wildly varies, furthermore, on how to deal with these types of words.

\item \textbf{[250]}, \textit{diachronic identity ... of the first} --- I.e., synchronic considerations still reign supreme. 

\end{itemize}

\subsection{Appendices}

\begin{itemize}

\item \textbf{[251]}, \textit{grammarian ... common element} --- This is an important claim if linguistics is to remain a descriptive, rather than prescriptive, science. Plato would reject the ``mistaken'' form, while Sextus Empiricus would accept only the ``mistaken'' form. Saussure finds use in both.

\item \textbf{[255]}, \textit{irreducible element ... root} --- Roots are still a very important morphemic category in contemporary linguistics.

\item \textbf{[258]}, \textit{it comes about ... forms compared} --- This is a claim that's fairly idiosyncratic to Indo-European languages. 

\end{itemize}

\newpage
\part{Notes on Chomsky's \textit{Syntactic Structures}}

\section{Preface/Introduction}

\begin{itemize}

\item \textbf{5}, \textit{the search for rigorous ... linguistic analysis} --- A summary statement of Chomsky's project. His rejection of ``logical niceties'' is a rejection of certain strands of analytical/mathematical philosophy then in vogue.

\item \textbf{5}, \textit{by pushing ... deeper understanding} --- Basically a statement of the scientific method.

\item \textbf{5}, \textit{linguistics who have questions} --- It's not exactly clear who he's taking about here. It's likely to refer either to a certain American strand of linguistics called the ``Structuralism'', of which Saussure was among the principle figure (taken up in a separate, European branch) or to the then-dominant trend in psychological theory called ``Behaviorism''.

\item \textbf{6}, \textit{three models} --- Again, these will be formal and falsifiable, hence Chomsky's scientific project.

\item \textbf{6}, \textit{W.V. Quine} --- 
This is one of the analytical philosophers Chomsky is probably arguing against.

\item \textbf{7}, \textit{This work was supported by ...} --- Chomsky's relationship to the U.S. military is a bit murky, especially since Chomsky was a vocal opponent of the Vietnam war a decade after this book came out. Some work about the influence of the military on his thought have come out relatively recently but that don't really have a good grasp of linguistic theory

\item \textbf{11}, \textit{by which sentences are constructed} --- This is a subtle shift from prior ways of doing syntax, from the process of description to that of production (i.e., something the mind can do).

\item \textbf{11}, \textit{a device} --- Also a departure from previous accounts. Again, syntax was priorly descriptive of certain strings in a language, but now it's an analyzable/delimitable/productive device.

\item \textbf{11}, \textit{with no specific reference to particular languages} --- The idea will be that differences in different languages are superficial. Chomsky's formulation of this problem is novel, but it taps into the same questions as Plato, Sextus Empiricus, and Saussure, all of whom try to figure out the general purpose/function of language. Hence, Plato's note that foreign languages have different words for the ``same'' thing, or Saussure's arbitrary nature of the sign.

\item \textbf{11}, \textit{representing utterances} --- This is a controversial statement. Some groups might say that this sets up a false paradigm for language, that it's not ``representable'' in some a priori ways. They and others might also try to reject such a procedure because of its association with Plato's representational view of the world.

\item \textbf{11}, \textit{simple and revealing} --- Important, if vague, criteria for Chomsky's scientific outlook.
 
\end{itemize}

\section{The Independence of Grammar}

\begin{itemize}

\item \textbf{13}, \textit{a set of sentences} --- This directly invokes innovations and terminology from the mathematics and set theory that started developing at the end of the 19th century.

\item \textbf{13}, \textit{separate the grammatical ... from the ungrammatical} --- Note that Chomsky, like Plato, Sextus, and Saussure, though in different ways, begins with a simple binary distinction. In practice, i.e., in contemporary linguistics, the notion of ``grammaticality'' is more complicated than this, though.

\item \textbf{13}, \textit{assume intuitive knowledge} --- This is a cornerstone of contemporary linguistic theory, based in the assumption that there is a language faculty in the human brain. Recall that Saussure methodologically rejected the need for native speaker intuition.

\item \textbf{14}, \textit{in many intermediate cases ... clear non-sentences} --- In practice what this means that, if some sentence is ambiguously grammatical, the theory can provide an explanation of why that's the case.

\item \textbf{14}, \textit{each grammar is related ... fixed in advance} --- I.e., Chomsky, like Saussure, wants the theory to include all (human) languages, but they differ as to how to methodologically access each language.

\item \textbf{15}, \textit{will project the finite ... grammatical utterances} --- I.e., it's possible to construct a number of theories to explain a pre-determined, finite set of strings, given enough rules. But this is not how language works.

\item \textbf{15}, \textit{a grammar mirrors the behavior of the speaker} --- Note again the huge shift here: syntax is about producing sentences. Note also the use of ``mirror'', to recall Plato's and Saussure's mimesis.

\item \textbf{15}, \textit{cannot be identified with ``meaningful''} --- This is another big move, that structure is a separate function/entity than ``meaning''. The role of meaning, the interaction between syntax and semantics, is still an important topic in contemporary linguistics.

\item \textbf{15}, \textit{Colorless green ideas sleep furiously} --- This is now a canonical example sentence. Note the method here, which is that example sentences are explicitly brought up, numbered and explained.

\item \textbf{16}, \textit{high order of statistical approximation} --- This is still a very controversial topic, which gets to the heart of linguistic theory. The question is: How specific does the brain's language module need to be? Namely, does the language module need in-built specific structures or can more general probabilistic functions account of language learning?

\item \textbf{17}, \textit{probabilistic modes give no particular insight} --- Probabilistic/computational models of language have come a long way since 1957, especially with regard to predicting unseen language data.

\end{itemize}

\section{An Elementary Linguistic Theory}

\begin{itemize}

\item \textbf{18}, \textit{levels of representations} --- Of course, linguistic theory has always made use of differing levels of phenomena for explanation. Chomsky makes the purpose of such a move more explicit here, namely the ability to make an explicit theory. It would be possible for the theory to be just a (near) one-to-one mapping from phenomena, but this would have no explanatory power.

\item \textbf{18}, \textit{one requirement ... that it be finite} --- This is the case because the brain is a finite machine.

\item \textbf{18}, \textit{suppose that we have a machine} --- These machines are usually called ``finite state automata'' (FSA). See \underline{\href{https://www.oxfordhandbooks.com/view/10.1093/oxfordhb/9780199276349.001.0001/oxfordhb-9780199276349-e-18}{this article}} for more information on FSA.

\item \textbf{20}, \textit{assign a probability} --- Probabilistic models similar to FSA are often called ``Markov models/chains''. As Chomsky notes in a footnote, this model is developed from mid-20th century information theory.

\item \textbf{20}, \textit{as the reader can easily convince himself} --- A somewhat typical dismissive tone from Chomsky

\item \textbf{21}, \textit{English is not a finite state language} --- This was a very important development for contemporary linguistic theory, though it's somewhat disputed.

\item \textbf{21}, \textit{each of these languages is not a finite state language} --- The proof, as Chomsky cites, is in \underline{\href{https://chomsky.info/wp-content/uploads/195609-.pdf}{this paper}} from a year earlier. The proof is a bit messy and also slightly disputed. The basic idea that certain types of languages can be formed, as in the (10) in the text, which include sentences that require an ever-expanding (i.e., non-finite) state machine.

\item \textbf{22}, \textit{all the mirror image properties} --- The formal truth of this is not exactly clear. But, again, the basic idea is that English includes the (infinitely) expanding set of sentences recursively built in this way, which would require a non-finite state machine.

\item \textbf{23}, \textit{arbitrarily decree ... for some fixed n} --- Note that this entails a sharp distinction between possible sentence and possible utterance. As Chomsky is assuming, the brain is finite, which means there must an upper limit as to the length of sentence that is producible (i.e., can be maintained in the memory), but, for Chomsky this is a separate issue than the potentially infinitely grammatical sentence.

\item \textbf{24}, \textit{if a grammar of this type ... will not produce} --- This is the main takeaway from the formal discussion.

\item \textbf{24}, \textit{at least one linguistic level cannot have this simple structure} --- Some work has indicated that phonology can be modelled on FSA.

\item \textbf{24}, \textit{generated from left to right} --- Note that this a primary concern. Contemporary linguistics, following Chomsky, will reject the linearity of language as part of its fundamental structure, making linearity a function of language's need to be spoken in time. Linearity, then, is less fundamental here than in Saussure.

\end{itemize}

\section{Phrase Structure}

\begin{itemize}

\item \textbf{26}, \textit{constituent analysis} --- As Chomsky notes, this is not a new concept. It refers to the idea that certain words are form definable phrases that are manipulable in unique ways. E.g., in the sentence ``he knew of the big dog'', ``knew of the big dog'' is a constituent, but ``he of the big dog'' is not.

\item \textbf{26}, \textit{new form of grammars} --- These and the examples given are called ``Context Free Grammars'' (CFG).

\item \textbf{26}, \textit{(13)} --- Note that the ``+'' symbol here and throughout doesn't mean a mathematical sum but instead refers to concatenation of strings (e.g., [``he'' + ``is'] means [``he is''].

\item \textbf{27}, \textit{represent the derivation} --- This is called a ``tree diagram'', and it is still a fundamental part of contemporary linguistic theory, though Chomsky is not the originator of it.

\item \textbf{28}, \textit{two derivations are equivalent} --- I.e., they produce the same literal string of words with the same constituencies.

\item \textbf{28}, \textit{constructional homonymity} --- This refers to a very common ambiguity between distinct interpretations. E.g., in the sentence ``the boy hit the bear with a stick'', ``with a stick'' can refer to the property of ``the boy'' or of ``the bear''. The hypothesis that these two different interpretations have different underlying structures.

\item \textbf{29}, \textit{only a single rule} --- This makes it so the system can't ``lose'' data by compressing two symbols into one.

\item \textbf{30}, \textit{terminal language} --- The idea will be, then, to think of possible English sentences as terminal strings in a terminal language and to develop constituency rules that can derive those terminal strings and only those terminal strings.

\item \textbf{30}, \textit{theorem} --- That every finite state language is terminal is true by definition.

\item \textbf{31}, \textit{introduced a symbol Z} --- This is an important point and a fundamental methodological move in linguistic theory. IT suggests that there are structural/real/operant parts of language that aren't necessarily spoken out loud.

\item \textbf{32}, \textit{no way of ordering} --- The idea is that, in the rules of a CFG, the rewriting of one symbol does not necessarily mean the new symbol is ``simpler'' in the sense of being, as a symbol, no longer relevant to the rule set.

\item \textbf{32}, \textit{phonemic structure} --- The phonological model here was influential but has mostly superseded by a model called ``Optimality Theory''.

\end{itemize}

\section{Limitations of Phrase Structure Description}

\begin{itemize}

\item \textbf{34}, \textit{I do not know whether} --- Some computational-based proofs in the eighties show that CFGs are not powerful enough to model natural languages.

\item \textbf{34}, \textit{ad hoc} --- What this means here is that the list of rules becomes so large that it approaches the number of strings under consideration, i.e., that the model there are a large number of unique rules for unique strings.

\item \textbf{36}, \textit{one of the best criteria} --- Conjunction analysis is still a diagnostic uses in contemporary linguistics, but it is relatively weak compared to many others.

\item \textbf{36}, \textit{(26)} --- Note that, as Chomsky admits, the rule in (26) is a different type than the rules in the CFG, though the CFG rules are still relevant.

\item \textbf{38}, \textit{essential reference to two distinct sentences} --- This is important when dealing with two or more sentences that seems very apparently connected in some way. E.g., the relationship between active and passive forms of a verb: ``the bear ate the boy'' and ``the boy was eaten by the bear''.

\item \textbf{40}, \textit{sequence of tense} --- This usually refers to the fact that in English, and many languages, tense of a subordinate clause is dependent on the tense of the controlling clause, e.g., ``he said he wanted a new life'' and not ``he said he wants a new life''. But Chomsky might be referring her to the fact that in most language, including English, tense/aspect are strictly ordered. In English the order is: modals < perfect < progressive < passive < verb.

\item \textbf{40}, \textit{no way to express the inversion} --- Because of the requirement that the lefthand side be unary in the CFG rules.

\item \textit{41}, \textit{have .. en} --- I.e., the requirement, as (30) shows, is that, e.g., the ``en'' suffix is dependent on ``have'', though it shows up on ``been''.

\item \textbf{42}, \textit{if the following V is transitive} --- I.e., the verb takes an object.

\item \textbf{43}, \textit{(34)} --- The point here is that rule (34) succinctly captures what seems to be an intuitive fact about English. The methodology is a little murky, since, as he admits, Chomsky hasn't proven this is, strictly speaking, necessary. The boundaries of native speaker intuition are not fully transparent.

\item \textbf{grammatical transformation} --- As Chomsky admits elsewhere, the concept of a ``transformation'' is taken from his PhD advisor, the linguist Zellig Harris.

\item \textbf{44-45}, \textit{this rule could not ... , etc.} --- The implication is that certain ``phases'' of the syntactic module occur before others, namely, certain structural transformations occur before spell-out of particular words.

\item \textbf{46}, \textit{something like this : (35)} --- That the syntax has two parts, the ``deep'' and ``surface'' structures, corresponding to CFG and transformational rules, is no longer accepted. Instead the syntax is considered to have only one binary operation called ``Merge''.

\item \textbf{47}, \textit{necessary to provide phrase structure only for kernel sentences} --- It becomes an empirical debate under this model as when base CFG rules and when transformational rules should be invoked.

\item \textbf{48}, \textit{grammar does not tell us how to synthesize a specific utterance} --- This would more likely be the task of psycholinguistics and neurology.

\end{itemize}

\section{On the Goals of Linguistic Theory}

\begin{itemize}

\item \textbf{49}, \textit{our fundamental concern} --- This paragraph is an explicit formulation as to how Chomsky is trying to conduct linguistics according the standard scientific method.

\item \textbf{49-50}, \textit{every grammar will ... any particular grammar} --- These are two key terms. They later become descriptive and explanatory adequacy in his 1965 book.

\item \textbf{50}, \textit{there is also no circularity in this conception} --- This is a constant concern of linguistic and other scientific theories. It roughly maps onto what is sometimes called the ``hermeneutic circle'' which describes the relationship between parts and a whole in, especially, literary texts.

\item \textbf{52}, \textit{it is unreasonable ... procedure for grammars} --- Chomsky is relatively uninterested, e.g., in trying to develop AI-like linguistic applications.

\item \textbf{53}, \textit{set of formal properties of grammars} --- E.g., the differences between FSA and CFG.

\item \textbf{54}, \textit{may not tell us ... from a corpus} --- This implies that the metatheory that develops could potentially produce languages that don't actually exist; i.e., the language module doesn't produce all possible languages. This is a distinct question from whether a grammar of a given language produces only the possible strings in that language.

\item \textbf{55}, \textit{until we give ... simplicity employed} --- As far as I know, this notion of ``simplicity'' has been relatively little defined, even through contemporary linguistics.

\item \textbf{56}, \textit{when we find that simplification ... right track} --- I.e., the goal is to find explanatory generalizations.

\item \textbf{56}, \textit{ultimate aim ... grammar of a language} --- This is an important methodological point, which displays Chomsky's interest in the general, underlying nature of the language module in the brain, rather than idiosyncratic behavioral mechanisms.

\item \textbf{58}, \textit{avoid all such problems ... walk+past} --- This is another instance in which a longstanding intuition in syntax, since Plato, is given a particular formal definition. Chomsky's formulation makes explicit the relation between different forms of the same root word. See footnote 8, as well, in which is explicit the fact that the morpheme \textit{past} shows up in syntactic analysis. The \textit{past} morpheme is an instance of a null entity, one that's not pronounced, appearing in the syntactic structure. This is a very productive methodological move that continues in contemporary linguistics.

\item \textbf{59-60}, \textit{it is quite true ...on lower levels} --- Contemporary linguistics continues to be divided into several subfields, including those that show up here: phonology, morphology, syntax, semantics.

\end{itemize}

\section{Some Transformations in English}

[Useful for examples as to how the theory works; notes would be even more technical. Let me know if there are questions here.]

\section{The Explanatory Power of Linguistic Theory}

\begin{itemize}

\item \textbf{85}, \textit{which no on may have ever produced} --- This is important for Chomsky's general notion w/r/t the poverty of the stimulus and language learning.

\item \textbf{86-87}, \textit{constructional homonymity, [etc.]} --- Three different situations are being described here:
    \begin{itemize}
    \item structural similarity: two different and different sounding phonemic sequences correspond to, at some level, the same structure
    \item structural ambiguity: the same phonemic sequence corresponds to two different structures
    \item constructional homonymity: a sub case of structural ambiguity, the same sound can be parsed in two ways and each parsing corresponds to a different structure
    \end{itemize}

\item \textbf{89}, \textit{to account for such phrases as (112i)} --- How to deal with English \textit{-ing} forms is still highly debated.

\item \textbf{89}, \textit{but both ... kernel sentences} --- The idea is that ``they growl the lions'' and ``the flowers raise'' are not grammatical sentences, so they don't provide ambiguity in the \textit{-ing} forms.

\item \textit{91}, \textit{sentence types} --- This is a less important task than classifying the transformational rules / phenomena themselves.

\end{itemize}

\section{Syntax and Semantics}

\begin{itemize}

\item \textbf{93}, \textit{the actual use of the language} --- The content of this distinction is not particularly clear.

\item \textbf{93}, \textit{with no appeal to meaning} --- This methodological choice is among the most controversial. In many ways contemporary linguistics has rejected the strict division between syntax and semantics, the latter of which has been formalized in, as Chomsky here requests, relatively precise ways.

\item \textbf{95}, \textit{(117i) cannot be accepted ... tokens, not types} --- This essentially constitutes an attack on Saussure's concept of the differential nature of the sign. Instead of each word being definitionally unique, Chomsky allows for the (perhaps more intuitive) notion that different words can mean the same thing.

\item \textbf{97}, \textit{operational criterion  for phonemic distinctness} --- Note the use of an experimental method defined by the object in question and which discovers two distinct phonemes, regardless of their meaning.

\item \textbf{98}, \textit{There is one further difficulty ... supposed to clarify} --- Compare the discussion here with Saussure's discussion on ``identity'' in ``Chapter 3'' of ``Synchronic Linguistics''.

\item \textbf{99}, \textit{[]footnote 7]} --- This seems to be another rejection of Saussure. Namely, utterances have positive meaning, but this is irrelevant to syntax. Utterances are also differentiated, which is significant to syntax, but this is not meaning. Saussure would say that meaning is essentially differential.

\item \textbf{100}, \textit{dummy carrier} --- Just as there are contentful but not-pronounced morphemes in the structure, so also are there pronounced but not-contentful morphemes.

\item \textbf{100}, \textit{meaning of some sort to such non-morphemes} --- This is what Socrates did in his (faux)-etymological musings in \textit{Cratylus}

\item \textbf{101}, \textit{[footnote 9} --- A somewhat less harsh statement of Chomsky's own positions.

\item \textbf{101}, \textit{might be missed} --- Note that this is a pragmatic concern, not a formal or theoretical one.

\item \textbf{102}, \textit{we can judge ... used and understood} --- Again, the notion of ``meaning'' compared to semantic ``use'' is not entirely clear.

\item \textbf{105}, \textit{sharply distinguish `grammatical' morphemes} --- Chomsky and contemporary theory ultimately do distinguish between lexical and functional morphemes like these.

\end{itemize}

\section{Summary}

[Nothing new here.]

\end{document}
